% =====================================================================
%  Arquitectura del repositorio Deasy (PUCE Sede Esmeraldas)
%  Documento de lectura. Compilar con:  pdflatex -interaction=nonstopmode
%  (dos pasadas para el indice)
% =====================================================================
\documentclass[11pt,a4paper,oneside]{report}

\usepackage[utf8]{inputenc}
\usepackage[T1]{fontenc}
\usepackage[spanish,es-tabla,es-noshorthands]{babel}
\usepackage{lmodern}
\usepackage[margin=2.4cm]{geometry}

\usepackage{array}
\usepackage{booktabs}
\usepackage{tabularx}
\usepackage{longtable}
\usepackage{enumitem}
\usepackage{xcolor}
\usepackage{framed}
\usepackage{listings}
\usepackage{fancyhdr}
\usepackage{parskip}
\usepackage[hidelinks]{hyperref}

% ---------------------------------------------------------------------
% Columnas para tablas
% ---------------------------------------------------------------------
\newcolumntype{L}{>{\raggedright\arraybackslash}X}
\newcolumntype{P}[1]{>{\raggedright\arraybackslash}p{#1}}

% ---------------------------------------------------------------------
% Cajas de aviso / concepto / nota
% ---------------------------------------------------------------------
\newenvironment{concepto}[1][Concepto para juniors]{%
  \par\medskip
  \colorlet{shadecolor}{blue!7}%
  \begin{shaded}\noindent\textbf{\color{blue!55!black}#1.}\ \ignorespaces}%
  {\end{shaded}\par\medskip}

\newenvironment{aviso}[1][Aviso]{%
  \par\medskip
  \colorlet{shadecolor}{orange!14}%
  \begin{shaded}\noindent\textbf{\color{orange!45!black}#1.}\ \ignorespaces}%
  {\end{shaded}\par\medskip}

\newenvironment{nota}[1][Nota]{%
  \par\medskip
  \colorlet{shadecolor}{black!6}%
  \begin{shaded}\noindent\textbf{#1.}\ \ignorespaces}%
  {\end{shaded}\par\medskip}

% ---------------------------------------------------------------------
% Codigo y diagramas ASCII
% ---------------------------------------------------------------------
\definecolor{codebg}{rgb}{0.97,0.97,0.97}
\definecolor{codekw}{rgb}{0.35,0.10,0.55}
\definecolor{codecm}{rgb}{0.30,0.45,0.30}
\definecolor{codestr}{rgb}{0.55,0.20,0.10}

\lstset{
  basicstyle=\ttfamily\footnotesize,
  backgroundcolor=\color{codebg},
  keywordstyle=\color{codekw}\bfseries,
  commentstyle=\color{codecm}\itshape,
  stringstyle=\color{codestr},
  breaklines=true,
  breakatwhitespace=false,
  columns=fullflexible,
  keepspaces=true,
  showstringspaces=false,
  frame=single,
  framerule=0.3pt,
  rulecolor=\color{black!25},
  xleftmargin=6pt,
  xrightmargin=2pt,
  aboveskip=8pt,
  belowskip=8pt,
  tabsize=2,
}

% Entorno para diagramas ASCII (sin resaltado, sin ajuste de linea)
\lstnewenvironment{diagrama}[1][]
  {\lstset{basicstyle=\ttfamily\scriptsize,language={},breaklines=false,#1}}
  {}

% Atajo para identificadores en linea
\newcommand{\rt}[1]{\texttt{#1}}

% ---------------------------------------------------------------------
% Cabeceras
% ---------------------------------------------------------------------
\pagestyle{fancy}
\fancyhf{}
\fancyhead[L]{\small\nouppercase{\leftmark}}
\fancyhead[R]{\small\thepage}
\renewcommand{\headrulewidth}{0.3pt}

\setcounter{tocdepth}{2}
\setcounter{secnumdepth}{3}

\title{\Huge\bfseries Arquitectura del sistema \textsc{Deasy}\\[0.4cm]
       \large Guia de lectura completa del repositorio}
\author{PUCE Sede Esmeraldas}
\date{9 de agosto de 2026}

\begin{document}

\pagenumbering{roman}
\maketitle

\begin{center}
\begin{minipage}{0.85\textwidth}
\small
\textbf{Sobre este documento.} Es una explicacion detallada y didactica de la arquitectura
del monorepo \textsc{Deasy}, escrita asumiendo cero contexto previo. Recorre el negocio,
las ocho piezas que corren en contenedores, el camino completo de una peticion HTTP, el
backend capa por capa, el modelo de datos y el motor de procesos, el frontend, el
microservicio de firma, la infraestructura y el despliegue, la estrategia de pruebas, las
trampas conocidas del repositorio y una ruta de lectura sugerida.

Todo lo que se afirma aqui procede de la lectura directa del codigo: rutas de fichero,
nombres de tabla, numeros de linea y conteos son reales en el momento de escribirlo.
\end{minipage}
\end{center}

\vfill
\tableofcontents
\clearpage
\pagenumbering{arabic}

% =====================================================================
\chapter{Que es Deasy (el problema de negocio)}
% =====================================================================

Antes de la arquitectura, el \textbf{para que}. Deasy es una plataforma de la
\textbf{PUCE Sede Esmeraldas} para gestionar los \textbf{tramites documentales recurrentes}
de la universidad: informes de gestion docente, requerimientos, informes de investigacion
formativa, memorandos, etc.

Hoy, sin Deasy, esos tramites viven en Word y PDF por correo: nadie sabe quien debe
entregar que, en que periodo, ni quien falta por firmar. Deasy modela cuatro cosas:

\begin{center}
\begin{tabularx}{\textwidth}{@{}P{3.2cm}L@{}}
\toprule
\textbf{Pregunta} & \textbf{Como la modela} \\
\midrule
?`Quien?           & Organigrama: unidades $\rightarrow$ puestos $\rightarrow$ personas ocupandolos \\
?`Que?             & Plantillas de documento, versionadas \\
?`Cuando?          & Periodos academicos (\rt{terms}) \\
?`Con que recorrido? & Flujos de \emph{entrega} (quien llena y revisa) y de \emph{firma} (quien firma y en que orden) \\
\bottomrule
\end{tabularx}
\end{center}

Y encima de todo eso, \textbf{firma electronica real} (PAdES, con certificados \rt{.p12} de
las autoridades certificadoras ecuatorianas: Security Data, ICERT y Firmasegura).

% =====================================================================
\chapter{Vista de pajaro: ocho piezas en contenedores}
% =====================================================================

Todo el sistema esta \textbf{dockerizado}. No hay un ``servidor'' monolitico, sino ocho
procesos que se hablan por red.

\begin{diagrama}
                        Navegador
                            | HTTPS
                            v
                  +--------------------+
                  |   nginx (proxy)    |  <- el portero: TLS + reparto de trafico
                  +----------+---------+
                     /                \
              "/"   /                  \  "/api/..."
                   v                    v
          +--------------+     +------------------+
          |   frontend   |     |     backend      |  Node 25 + Express 5
          |    Vue 3     |     |    (API REST)    |  + Socket.IO (tiempo real)
          +--------------+     +--+------+------+-+
                                  |      |      |
             +--------------------+      |      +--------------------+
             v                           v                           v
    +-----------------+       +-------------------+       +-----------------+
    |  PostgreSQL 17  |       |       MinIO       |       |    RabbitMQ     |
    |  (los datos)    |       |  (los ficheros)   |       |   (la cola)     |
    +-----------------+       +---------+---------+       +--------+--------+
                                        |                          |
                                        +------------+-------------+
                                                     v
                                          +----------------------+
                                          |        signer        |  Python + pyHanko
                                          |     (firma PDFs)     |  + helper Node "sigmaker"
                                          +----------------------+
\end{diagrama}

Traduccion de cada pieza para alguien que empieza:

\begin{itemize}[leftmargin=1.4em]
  \item \textbf{nginx}: el ``portero'' del edificio. Es lo unico expuesto a internet.
        Termina el HTTPS (descifra) y, segun la URL, manda la peticion al frontend o al
        backend. Tambien es lo que permite que el frontend llame a \rt{/api/...} sin
        problemas de CORS.
  \item \textbf{frontend}: la aplicacion web (Vue 3). En produccion es HTML, JS y CSS
        estatico servido por un nginx pequenito dentro de su propia imagen.
  \item \textbf{backend}: la API. Recibe peticiones HTTP, decide que se puede hacer, guarda
        en PostgreSQL, sube ficheros a MinIO y encola trabajos de firma en RabbitMQ.
  \item \textbf{PostgreSQL}: la base de datos relacional. \textbf{Es el unico motor de
        datos}: MariaDB y MongoDB fueron retirados. Veras comentarios ``ex-MongoDB'' en el
        codigo; son cicatrices de esa migracion.
  \item \textbf{MinIO}: almacenamiento de objetos compatible con S3. Los PDFs, plantillas,
        fotos de perfil y certificados \textbf{no} se guardan en la base de datos ni en el
        disco del backend: van aqui, en \emph{buckets} (contenedores de primer nivel).
  \item \textbf{RabbitMQ}: una cola de mensajes. Sirve para que el backend le pida al signer
        ``firma este PDF'' \textbf{sin quedarse bloqueado} esperando.
  \item \textbf{signer}: microservicio Python que hace la firma criptografica real.
  \item \textbf{analytics}: hoy es un marcador de posicion (\rt{CMD ["sleep","infinity"]}),
        no hace nada.
\end{itemize}

\begin{concepto}[Por que esa separacion importa]
Si la firma se hiciera dentro del backend, un PDF de 200 paginas bloquearia el servidor para
todos los usuarios mientras se procesa. Al ponerlo detras de una cola, el backend delega el
trabajo pesado y sigue atendiendo peticiones normalmente.
\end{concepto}

% =====================================================================
\chapter{El recorrido completo de una peticion}
% =====================================================================

Esto es lo que mas ayuda a orientarse en un repositorio nuevo. Sigue una peticion desde el
clic del usuario hasta la respuesta:

\begin{diagrama}
1. El navegador pide  https://localhost:8443/api/deasy/v1/users/me
                                             \___/\__________/
                                             nginx    backend

2. nginx (nginx/app-conf.d/default.conf.template):
   location /api/ { proxy_pass http://backend:3030/; }
                                                   ^
                           esa barra final BORRA el prefijo /api

3. El backend recibe   /deasy/v1/users/me
   (el prefijo /deasy/v1 se define en backend/config/apiPaths.js)

4. backend/index.js  ->  app.use(ROUTES.USERS, user_router)

5. backend/routes/user_router.js  ->  authMiddleware
                                  ->  loadAccessContext
                                  ->  requirePermissions("account.read")
                                  ->  controlador

6. backend/controllers/users/user_controler.js  ->  llama a UN servicio

7. backend/services/...  ->  SQL contra PostgreSQL via getPostgresPool()

8. Vuelta:  res.json({...})  ->  nginx  ->  navegador
\end{diagrama}

Ese es el noventa por ciento del sistema. Todo lo demas son variaciones sobre este mismo
esqueleto.

% =====================================================================
\chapter{El backend en detalle}
% =====================================================================

Es la pieza mas grande: unas 30.000 lineas de JavaScript \textbf{ESM} (\rt{import} y
\rt{export}, no \rt{require}), Express 5, sin ORM.

\section{Las capas y la regla no negociable del proyecto}

\begin{diagrama}
routes/       "que URL?"       Solo define paths y encadena middlewares
   |
   v
controllers/  "transporte"     Valida entrada, llama a UN servicio, responde HTTP
   |
   v
services/     "la logica"      Reglas de negocio, transacciones, maquinas de estado
   |
   v
config/postgres.js             Acceso a datos (pool de conexiones)
\end{diagrama}

El fichero \rt{CLAUDE.md} del repositorio lo dice sin rodeos: \textbf{``Los controllers son
transporte, no logica''}. Si un controller pasa de unas 40 lineas o abre una transaccion,
hay que extraer un servicio. Esto no siempre se ha respetado --- hay infractores historicos
como \rt{user\_controler.js} con 1.513 lineas --- y estan catalogados como deuda tecnica en
\rt{docs/planes/referencia/calidad-y-medicion.md}.

\subsection{Que hay dentro de cada capa}

\begin{itemize}[leftmargin=1.4em]
  \item \textbf{\rt{routes/}}: 14 ficheros planos (\rt{user\_router.js}, \rt{sign\_router.js},
        \rt{admin\_router.js}, \rt{chat\_router.js}, \rt{dossier\_router.js}, etc.). Todos
        cuelgan del prefijo \rt{/deasy/v1}.
  \item \textbf{\rt{controllers/}}: 8 subcarpetas por dominio (\rt{users/}, \rt{admin/},
        \rt{sign/}, \rt{chat/}, \rt{tareas/}, \rt{system/}, \rt{whatsapp/}, \rt{empresa/}),
        27 ficheros en total.
  \item \textbf{\rt{services/}}: 18 subcarpetas, unos 66 ficheros. Aqui estan los dominios
        reales: \rt{admin/} (con sub-subcarpetas \rt{kernel/}, \rt{crud/}, \rt{templates/},
        \rt{processes/}, \rt{org/}, \rt{generation/}), \rt{auth/}, \rt{chat/},
        \rt{documents/}, \rt{sign/}, \rt{system/}, \rt{tasks/}, \rt{users/}, \rt{mail/},
        \rt{storage/}, \rt{realtime/}, \rt{infrastructure/}, \rt{whatsapp/}, \rt{external/}.
  \item \textbf{\rt{database/}}: solo dos ficheros --- \rt{postgres\_initializer.js} y
        \rt{postgres\_schema.sql}. No es una capa de repositorios, es el arranque del esquema.
\end{itemize}

Un detalle de diseno elegante en \rt{services/admin/}: la carpeta \rt{kernel/} \textbf{no
importa a nadie} y todos importan de ella. Nada apunta ``hacia arriba''. Eso evita
dependencias circulares, que en JavaScript producen errores dificilisimos de diagnosticar en
tiempo de ejecucion.

\subsection{Inventario de routers}

\begin{center}
\footnotesize
\begin{tabularx}{\textwidth}{@{}P{4.2cm}P{3.1cm}P{1.3cm}L@{}}
\toprule
\textbf{Router} & \textbf{Montaje} & \textbf{Lineas} & \textbf{Autenticacion} \\
\midrule
\rt{user\_router.js}           & \rt{/users}         & 318 & mixta (login y create son publicos) \\
\rt{admin\_router.js}          & \rt{/admin}         & 32  & \rt{authMiddleware} + \rt{loadAccessContext} a nivel de router \\
\rt{sql\_admin\_router.js}     & \rt{/admin/sql}     & 139 & hereda del padre + \rt{requireSqlAdminPermission} \\
\rt{sign\_router.js}           & \rt{/sign}          & 93  & por ruta \\
\rt{dossier\_router.js}        & \rt{/dossier}       & 85  & \rt{authMiddleware} + \rt{loadAccessContext} \\
\rt{chat\_router.js}           & \rt{/chat}          & 53  & \rt{authMiddleware} (sin RBAC) \\
\rt{whatsapp\_router.js}       & \rt{/whatsapp}      & 25  & \textbf{ninguna} \\
\rt{tarea\_router.js}          & \rt{/tarea}         & 15  & solo en \rt{/supervised-stuck} \\
\rt{reset\_password\_router.js}& \rt{/reset-password}& 14  & publica \\
\rt{notification\_router.js}   & \rt{/notifications} & 12  & \rt{authMiddleware} \\
\rt{email\_router.js}          & \rt{/email}         & 9   & publica \\
\rt{program\_router.js}        & \rt{/program}       & 9   & \textbf{ninguna} \\
\rt{unit\_router.js}           & \rt{/units}         & 9   & \textbf{ninguna} \\
\rt{system\_router.js}         & \rt{/system}        & 9   & publica (bootstrap) \\
\bottomrule
\end{tabularx}
\end{center}

\begin{nota}[Duplicaciones reales]
\rt{program\_router.js} y \rt{unit\_router.js} son \emph{identicos}: ambos exponen
\rt{getPrograms} y \rt{createProgram} del mismo controlador. Y \rt{notification\_router.js}
duplica dos endpoints que ya estan en \rt{chat\_router.js}.
\end{nota}

\section{El punto de entrada: \rt{backend/index.js}}

Monta los middlewares en este orden exacto:

\begin{center}
\footnotesize
\begin{tabularx}{\textwidth}{@{}P{0.8cm}P{5.6cm}L@{}}
\toprule
\textbf{\#} & \textbf{Que monta} & \textbf{Detalle} \\
\midrule
0 & \rt{import "dotenv/config"} & primera linea del fichero \\
1 & \rt{app.set("trust proxy", 1)} & va detras de nginx \\
2 & \rt{app.use(cors(...))} & envuelto en un middleware propio; loguea el origen en \emph{cada} peticion \\
3 & \rt{app.use(express.json())} & parser JSON global \\
4 & \rt{app.use(cookieParser())} & necesario para la cookie \rt{refreshToken} \\
5 & \rt{app.get("/health")} & llama a \rt{assertPostgresConnection()}; 200 o 503 \\
6 & Swagger UI & en \rt{/deasy/docs} \\
7 & \rt{/deasy/docs.json} & el JSON crudo del spec \\
8 & Los 14 routers & ver tabla anterior \\
9 & \rt{express.static("public")} & la carpeta \rt{backend/public/} esta vacia \\
\bottomrule
\end{tabularx}
\end{center}

El arranque (\rt{startServer()}) hace, en secuencia:

\begin{diagrama}
startServer()
 |- initializeDatabaseWithRetry()      // bucle de reintentos
 |    |- assertPostgresConnection()    // config/postgres.js
 |    +- ensurePostgresSchema({reset}) // database/postgres_initializer.js
 |- http.createServer(app)
 |- realtimeGateway.init(httpServer, { corsOrigin, credentials: true })
 +- httpServer.listen(PORT)            // PORT || 3030
\end{diagrama}

Los reintentos se parametrizan con \rt{DB\_INIT\_MAX\_ATTEMPTS} (20),
\rt{DB\_INIT\_RETRY\_DELAY\_MS} (3000) y \rt{DB\_RESET\_SCHEMA\_ON\_START}. Si se agotan los
reintentos, \textbf{el backend no aborta}: loguea el fallo y sigue escuchando.

\begin{nota}[Que NO se conecta al arrancar]
RabbitMQ no abre conexion (se usa su API HTTP bajo demanda). MinIO construye el cliente al
importar el modulo, pero no crea buckets al arrancar. El bot de WhatsApp solo se inicializa
por peticion explicita. Y Swagger se configura con \rt{apis: []}, o sea que \textbf{no}
escanea anotaciones JSDoc: todo el spec es un objeto estatico en
\rt{backend/config/swagger/definition.js} (581 lineas).
\end{nota}

\section{El acceso a datos: un adaptador peculiar}

Aqui hay algo que confunde si no se explica. El fichero \rt{backend/config/postgres.js}
(512 lineas) \textbf{no es un cliente de PostgreSQL normal}. Es un \textbf{adaptador que
finge ser \rt{mysql2/promise}}.

?`Por que? El proyecto \textbf{migro de MariaDB a PostgreSQL}. Habia unos 881 sitios en el
codigo escribiendo SQL en estilo MySQL. Reescribirlos todos era inviable, asi que se escribio
una capa que traduce al vuelo. Hace tres cosas.

\subsection{\rt{bindParams}: traduce los marcadores de posicion}

MySQL usa \rt{?}; PostgreSQL usa \rt{\$1}, \rt{\$2}:

\begin{lstlisting}[language=Java]
// Lo que escribe el programador:
pool.query("SELECT * FROM persons WHERE cedula = ? AND status = ?", [ced, "Activo"])

// Lo que llega a PostgreSQL:
"SELECT * FROM persons WHERE cedula = $1 AND status = $2"
\end{lstlisting}

Y es cuidadoso: no toca los \rt{?} que estan dentro de comentarios (\rt{--}, \rt{/* */}) o de
literales de texto, porque mantiene una tabla de ``tramos protegidos''. Ademas expande
parametros de tipo array igual que hacia mysql2: un escalar da \rt{\$n}; un array vacio da
\rt{NULL}; \rt{[1,2,3]} da \rt{\$1, \$2, \$3}; y \rt{[[1,2],[3,4]]} da \rt{(\$1,\$2), (\$3,\$4)}.

Un cambio reciente e importante: \textbf{falla ruidosamente si faltan parametros}. Antes
mandaba \rt{NULL} en silencio, que es peor porque el fallo aparece lejos de su causa. El
mensaje de error, a proposito, \emph{no} incluye el SQL: varios controllers responden
\rt{error.message} directamente al cliente.

\subsection{\rt{translateDialect}: traduce funciones}

\rt{IFNULL} $\rightarrow$ \rt{COALESCE}; \rt{GROUP\_CONCAT} $\rightarrow$ \rt{string\_agg};
\rt{CURDATE()} y \rt{CURTIME()} $\rightarrow$ sus equivalentes; \rt{FROM DUAL} desaparece;
\rt{<=>} $\rightarrow$ \rt{IS NOT DISTINCT FROM}; \rt{INSERT IGNORE} $\rightarrow$
\rt{ON CONFLICT DO NOTHING}; \rt{IF()} $\rightarrow$ \rt{CASE WHEN}; \rt{FIELD()}
$\rightarrow$ un \rt{CASE} posicional; \rt{YEAR()}, \rt{MONTH()} $\rightarrow$ \rt{EXTRACT}.

El caso mas delicado es \rt{ON DUPLICATE KEY UPDATE} $\rightarrow$
\rt{ON CONFLICT (...) DO UPDATE SET}, porque PostgreSQL exige que le digas \emph{sobre que
columnas} es el conflicto. El adaptador lo \textbf{infiere consultando el catalogo} de
PostgreSQL (\rt{pg\_index}) y cachea el resultado en memoria.

\subsection{Transacciones al estilo mysql2}

\begin{lstlisting}[language=Java]
const conn = await pool.getConnection();
await conn.beginTransaction();     // ejecuta BEGIN
try {
  await conn.query("INSERT ...");
  await conn.query("UPDATE ...");
  await conn.commit();
} catch (e) {
  await conn.rollback();
} finally {
  conn.release();                  // devuelve la conexion al pool
}
\end{lstlisting}

Hay 11 ficheros de produccion que usan transacciones, entre ellos
\rt{services/admin/generation/launch.js}, \rt{services/admin/templates/templateLifecycle.js},
\rt{services/documents/FillRequestWorkflowService.js}, \rt{services/sign/PdfSigningService.js}
y \rt{services/system/SystemBootstrapService.js}.

\begin{concepto}[Que es un ``pool'' de conexiones]
Abrir una conexion a la base de datos es caro (decenas de milisegundos: handshake de red,
autenticacion, negociacion). Un \emph{pool} mantiene N conexiones ya abiertas y te presta una
cuando la necesitas. Aqui \rt{max = 10}. Por eso es \textbf{critico} llamar siempre a
\rt{release()}: si te olvidas, se agotan las diez y la aplicacion entera se cuelga esperando
una conexion que nunca vuelve.
\end{concepto}

\begin{aviso}[El SQL no lo valida nadie hasta que se ejecuta esa rama]
Textual del \rt{CLAUDE.md}. El SQL es una cadena de texto: \rt{node --check} no la mira, el
detector de imports tampoco, y el backend arranca igual. Asi sobrevivieron meses \emph{cuatro}
sentencias \rt{UPDATE ... INNER JOIN ... SET} (sintaxis multi-tabla de MySQL que PostgreSQL
\textbf{rechaza}) que dejaron el endpoint \rt{POST /sign/fill-requests/:id/return}
\textbf{roto para todo el mundo}.

PostgreSQL quiere \rt{UPDATE tabla alias SET col = ... FROM otra WHERE union AND filtros}, con
las columnas del \rt{SET} \textbf{sin cualificar}. Y \rt{grep "UPDATE.*JOIN"} \emph{no
encuentra nada}, porque el SQL ocupa varias lineas.
\end{aviso}

\subsection{Donde vive el SQL}

\rt{getPostgresPool} se importa en \textbf{28 ficheros}: 17 en \rt{services/}, 5 en
\rt{controllers/}, y el resto en \rt{config/}, \rt{database/} y \rt{utils/}. Es decir,
\textbf{no hay una capa de repositorios unica} --- es una deuda conocida. Solo
\rt{services/auth/} tiene repositorios nominales (\rt{UserRepository.js},
\rt{UserCertificateRepository.js}).

\section{El CRUD generico de administracion}

Esta es una de las piezas mas ingeniosas del backend. En vez de escribir un controlador para
cada una de las 44 tablas administrables, hay tres ficheros que colaboran:

\begin{itemize}[leftmargin=1.4em]
  \item \textbf{\rt{backend/config/sqlTables.js}} (1.001 lineas): un \textbf{catalogo
        declarativo} de metadatos. Para cada tabla describe sus campos, etiquetas humanas,
        tipos, cuales son obligatorios y cuales de solo lectura.
  \item \textbf{\rt{backend/services/admin/SqlAdminService.js}} (901 lineas): lee ese catalogo
        y \textbf{construye los SELECT, INSERT, UPDATE y DELETE dinamicamente}.
  \item \textbf{\rt{backend/services/admin/crud/tableHooks.js}} (1.120 lineas): los
        ``injertos'' --- la logica especial que necesita cada tabla concreta.
\end{itemize}

Una entrada del catalogo tiene esta forma:

\begin{lstlisting}[language=Java]
{
  table: "unit_types",
  label: "Tipos de unidad",
  category: "Estructura",
  primaryKeys: ["id"],
  fields: [
    { name: "name", label: "Nombre", type: "text", required: true },
    ...
  ],
  searchFields: ["name"]
}
\end{lstlisting}

\textbf{Ventaja}: anadir una tabla al panel de administracion es anadir una entrada al
catalogo, no escribir un CRUD entero. \textbf{Desventaja}: \rt{tableHooks.js} se ha convertido
en un cajon de sastre. El \rt{CLAUDE.md} avisa: \emph{``No injertes casos especiales en el
camino generico. Es el olor que hizo God a AdminTableManager''}.

\section{Autenticacion y autorizacion}

Son \textbf{dos cosas distintas} y conviene no confundirlas nunca:

\begin{itemize}[leftmargin=1.4em]
  \item \textbf{Autenticacion} = ``?`quien eres?'' $\rightarrow$ JWT
  \item \textbf{Autorizacion} = ``?`puedes hacer esto?'' $\rightarrow$ RBAC
\end{itemize}

\subsection{Autenticacion (JWT)}

\begin{diagrama}
POST /users/login  { cedula, password }
   |
   v
AuthService.login()  ->  bcrypt.compare(password, password_hash)
   |
   v
Devuelve DOS tokens:
   * access token   -> 2 horas   -> en el JSON; el frontend lo guarda en localStorage
   * refresh token  -> 30 dias   -> en una COOKIE httpOnly (el JS no puede leerla)
\end{diagrama}

\begin{concepto}[Por que dos tokens]
El \emph{access token} viaja en cada peticion, asi que si te lo roban el dano dura poco (dos
horas). El \emph{refresh token} solo se usa para pedir uno nuevo, y al ser una cookie
\rt{httpOnly}, un script malicioso inyectado en la pagina \textbf{no puede leerlo} desde
JavaScript. Es defensa en profundidad: dos secretos con exposiciones distintas.
\end{concepto}

Detalles concretos: el access token se firma con \rt{JWT\_SECRET} y expira en
\rt{60*60*2} segundos; el refresh con \rt{JWT\_REFRESH} y expira en \rt{60*60*24*30}. La
cookie lleva \rt{httpOnly: true} y \rt{secure} activo salvo en modo desarrollador.

El payload del JWT es \textbf{solo \rt{\{ uid \}}}. Nada de roles ni permisos dentro. Eso es
deliberado: si metieras los permisos en el token, revocarle un rol a alguien no tendria efecto
hasta que expirase. Aqui los permisos se resuelven \textbf{contra la base de datos en cada
peticion}.

La politica de contrasenas vive en \rt{backend/utils/passwordPolicy.js}
(\rt{evaluatePasswordPolicy}, exige 3 de 5 criterios) y la aplica
\rt{backend/middlewares/val\_password.js}, que ademas \emph{hashea en el sitio}
\rt{req.body.password} con bcrypt y salt 10 antes de llamar a \rt{next()}.

\subsection{Autorizacion (RBAC)}

RBAC son las siglas de \emph{Role-Based Access Control}, control de acceso basado en roles.
El modelo es:

\begin{diagrama}
resources x actions  -->  permissions ("dossier.read", "signature_flows.update", ...)
                                |
                                v
                          role_permissions  -->  roles
                                                   |
                                                   v
                                            role_assignments (persona + rol + UNIDAD)
                                                   ^ derived_from_assignment_id
                                            position_assignments  <--  cargo_role_map
\end{diagrama}

Trece recursos (\rt{account}, \rt{dossier}, \rt{security}, \rt{people}, \rt{units},
\rt{academic\_terms}, \rt{process\_definitions}, \rt{process\_execution}, \rt{templates},
\rt{documents}, \rt{fill\_flows}, \rt{signature\_flows}, \rt{contracts}) por cinco acciones
(\rt{read}, \rt{create}, \rt{update}, \rt{delete}, \rt{manage}) dan \textbf{65 permisos}. Y hay
\textbf{13 roles}.

Dos sutilezas importantes:

\begin{enumerate}[leftmargin=1.6em]
  \item \textbf{Los roles son contextuales a una unidad.} Puedes ser \rt{GestorProcesos} en la
        Facultad de Ingenieria y nadie en el resto de la universidad. La tabla
        \rt{role\_assignments} lleva \rt{unit\_id} y un \rt{max\_depth}: hasta que profundidad
        del subarbol organizativo se hereda el rol.
  \item \textbf{Hay roles derivados automaticamente del puesto.} La tabla \rt{cargo\_role\_map}
        dice, por ejemplo, ``el cargo DOCENTE otorga el rol GestorEjecucionProcesos''. Y un
        \textbf{trigger de PostgreSQL} (\rt{trg\_position\_assignments\_after\_insert}) crea
        esas asignaciones con \rt{source='derived'} cuando alguien ocupa el puesto, y las
        revoca cuando lo deja. Ese mismo trigger \textbf{reasigna los entregables abiertos} al
        nuevo ocupante.
\end{enumerate}

Los middlewares estan en \rt{backend/middlewares/rbac.js}:

\begin{center}
\footnotesize
\begin{tabularx}{\textwidth}{@{}P{6.2cm}L@{}}
\toprule
\textbf{Middleware} & \textbf{Que comprueba} \\
\midrule
\rt{loadAccessContext}                & Carga roles y permisos en \rt{req.auth} / \rt{req.access}; 401 si el usuario no existe o esta inactivo \\
\rt{requirePermissions(reqs, \{all\})}& Que tenga el permiso (OR por defecto, AND con \rt{\{all:true\}}) \\
\rt{requireAnyRole(roles)}            & Que tenga alguno de los roles indicados \\
\rt{requireRouteUserAccess(\{...\})}  & Que sea \textbf{el dueno} del recurso \textbf{o} tenga rol elevado, \emph{y} ademas el permiso \\
\rt{requireCedulaAccess(\{...\})}     & Igual, pero comparando por cedula \\
\rt{requireDossierAccess(action)}     & Azucar sintactico sobre el anterior con \rt{resource: "dossier"} \\
\rt{requireSqlAdminPermission(...)}   & Deduce el recurso desde \rt{req.params.table} y la accion desde el metodo HTTP \\
\bottomrule
\end{tabularx}
\end{center}

\begin{aviso}[Por que existe \rt{requireDossierAccess}]
Por un \textbf{IDOR} real y cerrado. IDOR (\emph{Insecure Direct Object Reference}) es el fallo
donde cambiando un identificador en la URL accedes a datos que no son tuyos. Aqui el guard
miraba \emph{la tarea} en vez de \emph{el entregable}, y un docente podia descargar el
documento de otro.
\end{aviso}

El catalogo canonico esta en \rt{backend/config/rbacCatalog.js} (229 lineas) y es tambien la
\textbf{fuente de siembra}: \rt{SystemBootstrapService.js} lo importa para poblar roles y
permisos en la base de datos, borrando y reescribiendo \rt{role\_permissions} en cada arranque.
O sea que \textbf{el catalogo del codigo es la fuente de verdad}, no la base de datos.

\section{Manejo de errores}

Hay una decision de diseno explicita: \textbf{no hay un middleware de error global}. Cada
controller hace su propio \rt{catch}:

\begin{lstlisting}[language=Java]
catch (error) {
  res.status(error.statusCode ?? 500).json({ message: error.message });
}
\end{lstlisting}

\rt{backend/errors/HttpError.js} define la clase y las fabricas: \rt{notFound()} (404),
\rt{forbidden()} (403), \rt{conflict()} (409), \rt{badRequest()} (400). El comentario de
cabecera del fichero dice literalmente que \emph{``un error SIN statusCode sigue siendo un 500
--- y eso esta bien''}.

Aparte, \rt{backend/errors/sqlErrors.js} traduce los codigos SQLSTATE de PostgreSQL a mensajes
humanos: \rt{23505} (violacion de UNIQUE) da ``ya existe una persona con esa cedula''. Y lo
hace comparando por \textbf{nombre de constraint} (\rt{error.constraint}), no buscando
subcadenas en el mensaje, que es fragil y se rompe al cambiar de idioma o de version.

Un cuarto elemento: \rt{backend/middlewares/uploadError.js} traduce los siete codigos de error
de \rt{multer} a mensajes en espanol y \textbf{oculta el mensaje si es un 5xx}. Se monta con
\rt{router.use(handleUploadError)} \emph{despues} de las rutas, en tres routers
(\rt{user\_router}, \rt{dossier\_router}, \rt{sign\_router}). Existe porque sin el, Express
devolvia HTML con el stack trace completo, incluidas las rutas absolutas dentro del contenedor.

\section{Integraciones externas}

\subsection{RabbitMQ, con sorpresa}

Aunque \rt{amqplib} figura en \rt{package.json}, \textbf{no se usa en ningun sitio}. Toda la
comunicacion va por la \textbf{API HTTP de gestion} de RabbitMQ
(\rt{backend/services/infrastructure/rabbitmq\_http.js}). El patron es \textbf{RPC con
polling}:

\begin{diagrama}
Backend genera correlationId = randomUUID()
  -> crea la cola de respuesta   deasy.sign.response.<uuid>
  -> publica el trabajo en       deasy.sign.request
  -> hace polling cada SIGN_POLL_MS (1000 ms) de la cola de respuesta
     hasta SIGN_TIMEOUT_MS (120000 ms)
\end{diagrama}

Las tres colas: \rt{deasy.sign.request}, \rt{deasy.sign.validate.request} y el prefijo
\rt{deasy.sign.response}. No hay consumidores persistentes en el backend: el consumidor real es
el microservicio signer.

\subsection{MinIO: siete buckets}

\begin{center}
\footnotesize
\begin{tabularx}{\textwidth}{@{}P{4.0cm}P{5.2cm}L@{}}
\toprule
\textbf{Bucket} & \textbf{Variable de entorno} & \textbf{Contenido} \\
\midrule
\rt{deasy-templates}    & \rt{MINIO\_TEMPLATES\_BUCKET}    & plantillas y semillas \\
\rt{deasy-documents}    & \rt{MINIO\_DOCUMENTS\_BUCKET}    & documentos generados \\
\rt{deasy-dossier}      & \rt{MINIO\_DOSSIER\_BUCKET}      & respaldos del CV personal \\
\rt{deasy-spool}        & \rt{MINIO\_SPOOL\_BUCKET}        & area temporal de firma \\
\rt{deasy-chat}         & \rt{MINIO\_CHAT\_BUCKET}         & adjuntos del chat \\
\rt{deasy-users}        & \rt{MINIO\_USERS\_BUCKET}        & fotos de perfil \\
\rt{deasy-certificates} & \rt{MINIO\_CERTIFICATES\_BUCKET} & certificados \rt{.p12} \\
\bottomrule
\end{tabularx}
\end{center}

La lista canonica esta en \rt{backend/scripts/lib/reset\_targets.mjs}.

\subsection{Socket.IO (tiempo real)}

\rt{services/realtime/RealtimeGateway.js} (179 lineas) esta montado sobre el \textbf{mismo
servidor HTTP} que Express, en el path \rt{/socket.io}. Sustituyo a un broker MQTT (EMQX).

\begin{itemize}[leftmargin=1.4em]
  \item \textbf{Autenticacion}: con \textbf{el mismo JWT} que la API, leido de
        \rt{socket.handshake.auth.token}. Al cliente siempre se le dice ``Token invalido''; el
        motivo real va al log.
  \item \textbf{Rooms} (mapeo uno a uno de los antiguos topics MQTT): \rt{user:\{personId\}}
        (join automatico), \rt{conversation:\{id\}}, \rt{process:\{id\}}.
  \item \textbf{Eventos entrantes}: \rt{conversation:subscribe}, \rt{conversation:unsubscribe},
        \rt{process:subscribe}, \rt{process:unsubscribe}, todos con callback de confirmacion.
        La suscripcion \textbf{reutiliza la autorizacion REST}, no la reimplementa.
  \item \textbf{Eventos salientes}: \rt{chat.message.created} y \rt{chat.notification.created}.
\end{itemize}

\subsection{Correo, WhatsApp y servicios ecuatorianos}

\rt{backend/lib/mailer.js} son once lineas de \rt{nodemailer} con el host SMTP
\textbf{escrito a fuego} en el codigo. El bot de WhatsApp
(\rt{services/whatsapp/WhatsAppBot.js}, 272 lineas) usa \rt{whatsapp-web.js} sobre Puppeteer
headless y muestra el QR de vinculacion por terminal; se inicializa manualmente y sus seis
endpoints \textbf{no llevan autenticacion}. Y \rt{services/external/webservices\_ec.js} valida
cedulas y numeros ecuatorianos contra un servicio externo.

\section{Los scripts del backend}

\begin{center}
\footnotesize
\begin{tabularx}{\textwidth}{@{}P{5.0cm}P{0.9cm}L@{}}
\toprule
\textbf{Fichero} & \textbf{L.} & \textbf{Para que} \\
\midrule
\rt{reset.mjs} & 114 &
  CLI de reset con \textbf{objetivos positivos y explicitos}:
  \rt{node scripts/reset.mjs <db|storage> [--yes]}. Sin argumentos imprime el uso y sale
  \textbf{sin tocar nada}. Antes eran tres ficheros con borrado por defecto y opt-out
  (\rt{--keep-db}); se invirtio la logica. \\
\rt{lib/reset\_targets.mjs} & 155 &
  Libreria: \rt{resetPostgres()}, \rt{resetMinio()} (vacia los 7 buckets) y
  \rt{describeContents()} (inventario previo al borrado). \\
\rt{bootstrap\_admin\_recovery.mjs} & 46 &
  Recupera o recrea el administrador. \rt{npm run recover:admin}. \\
\rt{check\_missing\_imports.mjs} & 124 &
  \textbf{Detector de ``simbolo movido pero no importado''}. Ver el aviso de abajo. \\
\rt{seed\_dev\_rich.mjs} & 284 &
  Fixture rica de desarrollo: crea un segundo proceso en modo \rt{single} y lo lanza, porque
  con el proceso unico del bootstrap la pantalla \rt{/home} no se puede verificar. \\
\rt{generate\_demo\_certificates.mjs} & 220 &
  Genera certificados PKCS\#12 autofirmados con openssl y los sube a
  \rt{deasy-certificates}, porque la interfaz solo permite \emph{subir} un \rt{.p12}. \\
\bottomrule
\end{tabularx}
\end{center}

\begin{aviso}[\rt{npm run check:imports} es obligatorio tras mover codigo]
\rt{node --check} valida \textbf{sintaxis, no imports}. Un simbolo movido sin su \rt{import} es
sintaxis perfectamente valida, el modulo \textbf{carga sin quejarse}, y revienta en tiempo de
\textbf{llamada}. Asi estuvieron rotos tres semanas cuatro \rt{ReferenceError} en produccion
(\rt{createUnitWithParent}, \rt{getCargoCodeMap},
\rt{getNextStorageVersionForTemplateCode}, \rt{loadTemplateArtifactMetaDocument} --- esta
ultima ya retirada con el \rt{meta.yaml}, pero sigue nombrada en el script como caso de prueba
historico).

Comprobar que el backend arranca \textbf{no sustituye} a ejecutar este script. Esta en CI.
\end{aviso}

\begin{nota}[No hay migraciones incrementales]
No existe nada tipo Flyway o Knex. El esquema se aplica \textbf{entero} en cada arranque desde
\rt{backend/database/postgres\_schema.sql}, apoyandose en que es idempotente
(\rt{CREATE TABLE IF NOT EXISTS}, \rt{CREATE OR REPLACE FUNCTION}, siembras con
\rt{ON CONFLICT}).
\end{nota}

% =====================================================================
\chapter{La base de datos y el motor de procesos}
% =====================================================================

\textbf{67 tablas} y una vista, en un unico fichero:
\rt{backend/database/postgres\_schema.sql} (1.642 lineas).

\begin{aviso}[Consecuencia de no tener migraciones]
El esquema entero se reaplica en cada arranque. Es simple y funciona bien sin datos de
produccion, pero significa que \textbf{no puedes hacer un \rt{ALTER TABLE} evolutivo} de forma
comoda: cualquier cambio de columna hay que pensarlo como idempotente o pasa por un reset.
\end{aviso}

\section{Vocabulario del dominio}

\begin{center}
\footnotesize
\begin{tabularx}{\textwidth}{@{}P{2.9cm}LP{3.6cm}@{}}
\toprule
\textbf{Termino} & \textbf{Que es en Deasy} & \textbf{Donde vive} \\
\midrule
\textbf{Proceso} &
  El tramite institucional en abstracto (``Informe de Gestion Docente''). Es solo un nombre y
  una jerarquia; el comportamiento esta en sus \emph{configuraciones}. &
  \rt{processes} \\
\textbf{Configuracion} &
  La version vigente del proceso: a quien alcanza, en que periodos corre, que entregables
  produce. Se activa y se retira. &
  \rt{process\_definition\_versions} \\
\textbf{Serie} &
  El \emph{eje} por el que un proceso se declina: por tipo de unidad, por cargo, o
  \rt{default} (sin variacion). &
  \rt{process\_definition\_series} \\
\textbf{Regla} &
  Quien recibe el proceso: que unidades, que cargo o puesto, con que politica de reparto. &
  \rt{process\_target\_rules} \\
\textbf{Corrida (run)} &
  El acto de \emph{lanzar} la configuracion en un periodo. Genera las tareas. &
  \rt{process\_runs} \\
\textbf{Tarea} &
  La instancia del proceso para un ambito concreto (una unidad, un periodo). Es el contenedor. &
  \rt{tasks} \\
\textbf{Entregable} &
  Cada documento concreto a producir dentro de la tarea, con responsable, vencimiento y estado. &
  \rt{task\_items} \\
\textbf{Plantilla} &
  El \emph{molde} del documento: schema de campos + cuerpo (Jinja2/LaTeX u ofimatico) +
  formatos, versionado y almacenado en MinIO. &
  \rt{deliverables} + \rt{template\_artifacts} \\
\textbf{Flujo de entrega} &
  Cadena de pasos ``quien llena y aprueba el documento antes de firmarlo''. &
  \rt{fill\_flow\_*} / \rt{fill\_requests} \\
\textbf{Firma} &
  Firma electronica PAdES sobre el PDF, con certificado \rt{.p12} del firmante. &
  \rt{signature\_flow\_*} / \rt{document\_signatures} \\
\textbf{Dossier} &
  El \textbf{expediente o CV personal} (titulos, experiencia, publicaciones). \emph{No} es el
  expediente de un proceso. &
  \rt{dossiers} + \rt{dossier\_items} \\
\bottomrule
\end{tabularx}
\end{center}

\section{El motor de procesos: serie, regla y flujo}

Esta es \textbf{la parte central del dominio} y la que mas cuesta al principio. La frase
mnemotecnica del proyecto:

\begin{center}
\emph{la \textbf{serie} nombra el proceso, la \textbf{regla} reparte el proceso,
el \textbf{flujo} reparte los pasos.}
\end{center}

\begin{diagrama}
processes                          "Informe de Gestion Docente"  (solo un nombre)
   |
   +--> process_definition_versions      LA CONFIGURACION (draft / active / retired)
          |    ^ series_id -> process_definition_series   <- LA SERIE (el eje de variacion)
          |
          +--> process_target_rules              <- LA REGLA: a que unidades y cargos alcanza
          +--> process_definition_period_types   <- en que tipos de periodo corre
          +--> process_definition_templates      <- que documentos produce (+ item_mode)
          |
          +--> process_runs                      <- EL LANZAMIENTO en un periodo concreto
                 |
                 +--> tasks                      <- una por (definicion x periodo x unidad)
                        |
                        +--> task_items          <- LOS ENTREGABLES concretos
                               |
                               +--> documents -> document_versions
                                      |
                                      +--> document_fill_flows / fill_requests      (entrega)
                                      +--> signature_flow_instances / signature_requests (firma)
\end{diagrama}

\subsection{Tabla por tabla}

\paragraph{\rt{processes}.} Solo identidad: \rt{name}, \rt{slug} unico, \rt{parent\_id}
autorreferencial, \rt{is\_active}. \textbf{Sin comportamiento.} Se hizo \rt{DROP} de
\rt{unit\_id}, \rt{program\_id}, \rt{person\_id} y \rt{term\_id} por vestigiales.

\paragraph{\rt{process\_definition\_series} --- la serie.}
\rt{source\_type} puede ser \rt{unit\_type}, \rt{cargo} o \rt{default}. Es el \textbf{eje de
variacion}: el mismo proceso puede tener una configuracion distinta por tipo de unidad o por
cargo. La funcion \rt{buildProcessDefinitionVersionName} genera el nombre resultante:
\emph{``<Proceso> por <Serie>''}, salvo en \rt{default}, que usa el nombre del proceso a secas.

\paragraph{\rt{process\_definition\_versions} --- la configuracion.}
Tiene \rt{status} (\rt{draft}, \rt{active}, \rt{retired}) y vigencia
(\rt{effective\_from} / \rt{effective\_to}). Dos mecanismos de integridad:

\begin{itemize}[leftmargin=1.4em]
  \item Una \textbf{columna generada} \rt{active\_series\_flag} mas un indice UNIQUE sobre
        \rt{(process\_id, variation\_key, active\_series\_flag)} garantiza \textbf{a nivel de
        base de datos} que solo hay una version activa por linea.
  \item Un \textbf{trigger} (\rt{trg\_process\_definition\_versions\_before\_update}) impide
        activarla si no tiene al menos una regla activa \emph{y} un tipo de periodo, lanzando
        una excepcion con mensaje en espanol.
\end{itemize}

\begin{concepto}[Reglas de negocio dentro de la base de datos]
Es poco habitual y merece la pena entender el porque: si la regla ``no se puede activar sin
reglas de alcance'' viviera solo en el codigo, cualquier script, migracion o consulta manual
podria saltarsela. Al ponerla en un trigger, \textbf{es imposible dejar la base de datos en un
estado invalido}, venga la escritura de donde venga.
\end{concepto}

\paragraph{\rt{process\_target\_rules} --- la regla.}
Reparte el alcance con tres piezas: \rt{unit\_scope\_type} (\rt{unit\_exact},
\rt{unit\_subtree}, \rt{unit\_type}, \rt{all\_units}), el \rt{cargo\_id} o \rt{position\_id}
dentro de la unidad, y una \rt{recipient\_policy} (\rt{all\_matches}, \rt{one\_per\_unit},
\rt{exact\_position}). Ademas \rt{priority} y vigencia.

\paragraph{\rt{process\_definition\_templates} --- el paquete de entregables.}
Vincula configuracion con plantilla: \rt{process\_definition\_id} + \rt{template\_artifact\_id},
mas \rt{sort\_order} y, sobre todo, \rt{item\_mode}.

\paragraph{\rt{process\_runs} --- la corrida.}
\rt{run\_mode} (\rt{automatic} o \rt{manual}), \rt{source\_run\_id} autorreferencial para
relanzamientos, \rt{created\_by\_user\_id}, \rt{reason} y \rt{status}. Relanzar es \textbf{una
corrida nueva}, no sobrescribir la anterior.

\paragraph{\rt{tasks}.}
Con un UNIQUE sobre \rt{(process\_definition\_id, term\_id, normalized\_scope\_unit\_id)} ---
donde la ultima es una columna generada \rt{COALESCE(scope\_unit\_id, 0)} --- que garantiza
\textbf{idempotencia del lanzamiento}: lanzar dos veces no duplica tareas.

\section{\rt{item\_mode}: los tres modos de emision}

Vive en \rt{process\_definition\_templates.item\_mode}, es decir, \textbf{en el vinculo
plantilla-proceso}, no en la plantilla. Es un matiz importante. Define \emph{cuando} nace el
entregable y \emph{de donde sale su flujo}:

\begin{center}
\footnotesize
\begin{tabularx}{\textwidth}{@{}P{2.2cm}P{3.4cm}LP{3.4cm}@{}}
\toprule
\textbf{Modo} & \textbf{Flujo} & \textbf{Instanciacion} & \textbf{Ejemplo} \\
\midrule
\rt{single} &
  Predefinido en la plantilla (entrega y firma) &
  Una instancia automatica al lanzar, por cada responsable que resuelvan las reglas &
  Informe de Investigacion Formativa \\
\rt{replicated} &
  Predefinido &
  \textbf{Sin fan-out automatico}: el responsable crea N replicas etiquetadas que
  \textbf{heredan} el flujo del original &
  Un requerimiento por cada docente nuevo \\
\rt{routed} &
  \textbf{Ninguno} &
  El usuario define entrega y firma \textbf{al instanciar} (en tiempo de ejecucion) &
  Tareas ad-hoc \\
\bottomrule
\end{tabularx}
\end{center}

Truco de implementacion para distinguirlos en la base de datos: los flujos \emph{de plantilla}
tienen \rt{task\_item\_id IS NULL}; los flujos creados en tiempo de ejecucion tienen
\rt{task\_item\_id != NULL}. Los materializa \rt{materializeRuntimeFlowForTaskItem}.

Las replicas son \rt{task\_items} con \rt{origin\_kind='user\_added'} enlazadas por
\rt{source\_task\_item\_id} al original.

El \textbf{``Proceso por defecto''} (slug \rt{default}, sembrado por el bootstrap) es un
\rt{routed} comodin para tareas ad-hoc que no pertenecen a ningun proceso (por ejemplo, ``haz
el informe de este evento''). El \rt{CLAUDE.md} insiste: \textbf{no es ``memorandums''}.

\begin{nota}[Resolutores de flujo autorables]
Solo \rt{task\_assignee} (``Responsable del entregable'') y \rt{cargo\_in\_scope} (``Por
cargo'') en plantillas \emph{official}; \rt{specific\_person} se anade en \emph{ad\_hoc}.
\rt{document\_owner}, \rt{position} y \rt{manual\_pick} estan \textbf{retirados}: ya no son una
deprecacion blanda de la autoria web, sino que \textbf{salieron del \rt{CHECK}} de
\rt{fill\_flow\_steps.resolver\_type} y \rt{signature\_flow\_steps.resolver\_type}, que hoy solo
admite esos tres valores. El \rt{ALTER} valida las filas existentes, asi que una base con un
valor retirado \textbf{no arranca}.

El criterio que los mato, y que conviene tener presente al anadir cualquier cosa a estas tablas:
\textbf{lo que la web no autora, no existe}. Su unico productor era el \rt{meta.yaml}; retirado
el YAML, se quedaron sin quien los escribiera.
\end{nota}

\section{Plantillas: el modelo ``libro y ediciones''}

\begin{diagrama}
template_seeds       (semillas del catalogo: source_path, preview_path, seed_type)
      |
      v
deliverables         EL LIBRO    -> code UNIQUE, display_name, owner_process_id,
                                    template_scope (official | ad_hoc), owner_person_id
      |  1:N
      v
template_artifacts   LA EDICION  -> storage_version, lifecycle_state, base_object_prefix,
                                    available_formats, schema_object_key, content_hash,
                                    parent_version_id (autorreferencial = linaje)
\end{diagrama}

Con UNIQUE sobre \rt{(deliverable\_id, storage\_version)}. El esquema lleva comentarios
explicitos: la identidad, el proceso propietario, el scope, la semilla y la persona propietaria
viven en \rt{deliverables}; \rt{template\_artifacts} guarda \textbf{solo} el estado y el
almacenamiento de cada version.

\begin{aviso}[La palabra ``entregable'' significa dos cosas]
En \rt{sqlTables.js} la tabla \rt{task\_items} se etiqueta ``Entregables'' (la \emph{instancia}
a producir). En el esquema, \rt{deliverables} es ``el libro'': la \emph{identidad de la
plantilla}. Es una fuente real de confusion al leer el codigo; siempre hay que mirar el
contexto.
\end{aviso}

\subsection{El contenido: Jinja2 sobre LaTeX}

El cuerpo del documento es un \textbf{contrato Jinja2 + LaTeX} empaquetado en MinIO, no en la
base de datos. El bootstrap publica la semilla \rt{backend/services/system/seeds/informe-general}
--- \rt{schema.json}, \rt{defaults.yaml}, \rt{README.md} y el arbol \rt{src/} con
\rt{main.tex.j2} y \rt{make.sh} --- y valida el pipeline: render Jinja2 con
\rt{StrictUndefined} $\rightarrow$ \rt{pdflatex} $\rightarrow$ PDF.

\textbf{No hay \rt{meta.yaml}}: el flujo se autora en la base, no en un YAML, desde el sub-paso
8 del frente 0. Y \rt{data.yaml} \textbf{no es un fichero de la semilla}: es un objeto de MinIO
que el bootstrap escribe copiando \rt{defaults.yaml} al prefijo del artifact
(\rt{publishBaseSeedAssets}).

La \textbf{propiedad} se expresa por el prefijo dentro de MinIO, en \rt{base\_object\_prefix}:
\rt{System/...} para las plantillas curadas por el administrador, \rt{Users/\{cedula\}/...}
para las subidas por un gestor. El pipeline de maduracion documentado es:

\begin{center}
\rt{office} (Word/Excel/PDF subido por el gestor) $\longrightarrow$
\rt{jinja2} (curado por el administrador)
\end{center}

\begin{nota}[Deuda declarada]
Falta \emph{cablear} el render Jinja2 a PDF en tiempo de ejecucion: hoy el render vive en el
\rt{make.sh} de la propia semilla (utileria de linea de comandos), y el backend no lo invoca.
\end{nota}

\subsection{Ciclo de vida y versionado}

\rt{template\_artifacts.lifecycle\_state} admite \rt{draft}, \rt{published} y \rt{retired}, y
\textbf{por defecto \rt{draft}}: una edicion nace sin publicar y hay que publicarla
explicitamente. El versionado es \textbf{por linaje}: una fila nueva con
\rt{storage\_version} nuevo y \rt{parent\_version\_id} apuntando a la anterior.
\rt{templateLifecycle.js} retira la version publicada previa del mismo codigo al publicar, y
valida que haya al menos un paso de entrega. Ese ultimo requisito \textbf{se relaja para
\rt{routed}}, que por definicion no autora flujo.

\section{Firmas}

\begin{diagrama}
signature_flow_templates  ->  signature_flow_steps  ->  signature_flow_instances
                                                              |
                                                              v
                                          signature_requests  ->  document_signatures
\end{diagrama}

\rt{signature\_flow\_steps} soporta \textbf{multiples firmantes por paso} (columna
\rt{signers} de tipo JSONB) con quorum configurable via \rt{approval\_mode}:

\begin{itemize}[leftmargin=1.4em]
  \item \rt{and}: firman \textbf{todos} los del paso.
  \item \rt{or}: basta \textbf{cualquiera}.
  \item \rt{at\_least}: un minimo de N, indicado en \rt{required\_signers\_min}.
\end{itemize}

La columna \rt{anchor\_refs} (JSONB) \textbf{es un fosil}: no tiene productor ni consumidor. El
escritor la serializa siempre como \rt{[]} y el lector la devuelve tal cual; ningun formulario le
pone un valor y ningun render la mira. Quien decide \textbf{donde se dibuja la firma} es la
columna \rt{slot} del paso, que el cuerpo Jinja2 embebe como
\rt{\{\{ signatures.<slot>.token \}\}}. Su gemela en el lado de entrega es
\rt{fill\_flow\_steps.field\_refs}, con el mismo problema.

Los catalogos de estado se siembran en el propio esquema:
\rt{signature\_statuses} (\rt{firmado}, \rt{fallido}, \rt{invalido}, \rt{cancelado}) y
\rt{signature\_request\_statuses} (\rt{pendiente}, \rt{en\_progreso}, \rt{completado},
\rt{rechazado}, \rt{cancelado}).

\subsection{El flujo de firma en lote}

\begin{center}
\footnotesize
\begin{tabularx}{\textwidth}{@{}P{5.4cm}L@{}}
\toprule
\textbf{Endpoint} & \textbf{Funcion} \\
\midrule
\rt{POST /sign/}                       & Firma individual \\
\rt{POST /sign/validate}               & Validacion de un PDF ya firmado \\
\rt{POST /sign/batch}                  & \textbf{Retirado} --- devuelve error y redirige a \rt{/batch/start} \\
\rt{POST /sign/batch/start}            & Arranca el trabajo asincrono (hasta 30 PDFs) \\
\rt{GET /sign/batch/:jobId}            & Estado del trabajo \\
\rt{GET /sign/batch/:jobId/download}   & Descarga el ZIP con los resultados \\
\bottomrule
\end{tabularx}
\end{center}

El estado se persiste en \rt{signature\_batch\_jobs} (clave \rt{job\_id}, mas \rt{status},
\rt{total}, \rt{processed}, \rt{success\_count}, \rt{failed\_count} y \rt{results} JSONB), con
un comentario revelador en el esquema: \emph{``persistido para sobrevivir reinicios del backend
(antes vivia en un Map en memoria)''}.

Separacion de responsabilidades: \rt{BatchSigningService.js} gestiona el ciclo del trabajo y el
bucle en segundo plano, pero \textbf{no} decide donde se guarda el PDF ni habla con el firmante
--- eso es de \rt{PdfSigningService.js}, que llama documento a documento. El ZIP se empaqueta
con \rt{/usr/bin/zip} (ruta absoluta a proposito).

Ademas del PDF, cada documento del lote lleva metadatos de contexto
(\rt{signatureRequestId}, \rt{documentVersionId}, \rt{processName}, \rt{unitLabel},
\rt{termName}, \rt{stepName}), que es lo que permite persistir la evidencia del flujo de
trabajo y no limitarse a firmar ficheros sueltos.

\section{El dossier (ex-MongoDB)}

El dossier es el CV o expediente personal: titulos, experiencia, publicaciones. Antes vivia en
MongoDB; ahora son \textbf{dos tablas}:

\begin{itemize}[leftmargin=1.4em]
  \item \rt{dossiers}: la raiz, UNIQUE por \rt{person\_id}, con la \rt{cedula} desnormalizada e
        indexada.
  \item \rt{dossier\_items}: \rt{dossier\_id} (con borrado en cascada) + \rt{section} +
        \textbf{\rt{data} JSONB} + \rt{url\_documento}.
\end{itemize}

Se conservo la flexibilidad de documento usando JSONB. La justificacion esta en el comentario
del esquema: \emph{``documento heterogeneo (CV con secciones) accedido siempre como arbol
completo por cedula''}.

Las antiguas ``colecciones'' de Mongo son ahora valores de \rt{section}: \rt{titulos},
\rt{experiencia}, \rt{referencias}, \rt{formacion}, \rt{certificaciones}, \rt{articulos},
\rt{libros}, \rt{ponencias}, \rt{tesis}, \rt{proyectos}.

Por compatibilidad, los \rt{\_id} se exponen como \textbf{String} para preservar el contrato
que tenia Mongo, y los valores por defecto de cada seccion replican exactamente los del antiguo
esquema de Mongoose.

\section{Los demas dominios de tablas}

\subsection{Identidad y personas}
\rt{persons} (con \rt{cedula}, \rt{email} y \rt{token} unicos, \rt{password\_hash},
\rt{status}), \rt{person\_certificates} (los \rt{.p12} en MinIO, con \rt{is\_default}),
\rt{email\_verification\_codes}, \rt{password\_reset\_codes}.

\subsection{Organizacion, unidades y puestos}
\rt{unit\_types}, \rt{units}, \rt{relation\_unit\_types}, \rt{unit\_relations},
\rt{cargos}, \rt{unit\_positions}, \rt{position\_assignments}. Dos garantias elegantes por
columna generada:

\begin{itemize}[leftmargin=1.4em]
  \item \rt{unit\_positions.head\_flag} garantiza \textbf{un solo jefe por unidad}.
  \item \rt{position\_assignments.current\_flag} garantiza \textbf{un solo ocupante actual por
        puesto}.
\end{itemize}

Y una vista, \rt{unit\_org\_levels}, que con una CTE recursiva sobre \rt{unit\_relations}
calcula el nivel organizativo, la unidad raiz y las unidades de nivel 2 y 3 de cada unidad.

Como subdominio propio esta la \textbf{contratacion}: \rt{vacancies}, \rt{aplications},
\rt{offers}, \rt{contracts}, \rt{contract\_origins} (con herencia por tabla) y
\rt{vacancy\_visibility}.

\subsection{Entregables, tareas y documentos}
\rt{tasks}, \rt{task\_items}, \rt{task\_assignments}, \rt{task\_item\_handovers} (bitacora de
traspaso de responsable), \rt{documents} (1:1 con \rt{task\_items}), \rt{document\_versions}
(con \rt{version} de tipo \rt{DECIMAL(4,1)}), \rt{document\_attachments} y
\rt{document\_workflow\_observations} (observaciones, devoluciones y rechazos).

\subsection{Chat y notificaciones}
\rt{chat\_conversations} (con tipos \rt{direct}, \rt{group}, \rt{thread}, \rt{process\_thread},
\rt{unit}), \rt{chat\_conversation\_participants}, \rt{chat\_messages},
\rt{chat\_message\_attachments}, \rt{chat\_message\_reads} y \rt{chat\_notifications}.

\begin{nota}[Decision de diseno explicita]
En el dominio de chat, \rt{person\_id}, \rt{process\_id} y \rt{unit\_id} son \textbf{claves
foraneas logicas sin constraint}, ``para no acoplar''. Y \rt{last\_message\_id} tampoco lleva FK,
para evitar un ciclo de dependencias entre tablas.
\end{nota}

\subsection{Auditoria}
\textbf{No existe una tabla de auditoria general} (nada de \rt{audit\_log} o \rt{bitacora}). La
trazabilidad es \textbf{especifica por dominio}: \rt{task\_item\_handovers} para los traspasos,
\rt{document\_workflow\_observations} para devoluciones, \rt{role\_assignments.revoked\_at} y
\rt{revoked\_reason} para revocaciones, \rt{signature\_batch\_jobs.results} para el resultado
por fichero, y \rt{created\_at} / \rt{updated\_at} con trigger \rt{set\_updated\_at()} en la
mayoria de tablas.

% =====================================================================
\chapter{El frontend}
% =====================================================================

\textbf{Vue 3 + Vite 8 + vue-router 5 + Tailwind v4 + Vitest}, sin TypeScript. 216 ficheros
en \rt{src/}: 126 \rt{.vue}, 87 \rt{.js} y 3 \rt{.css}.

\section{Organizacion: modular por dominio}

\begin{diagrama}
src/
+-- core/       infraestructura transversal (router, httpClient, apiConfig, accessControl)
+-- layouts/    esqueletos de pagina (AppWorkspaceShell, headers, menus)
+-- modules/    LOS DOMINIOS
|   +-- admin/     (~80 ficheros)  panel de administracion + grafos
|   +-- home/      (26)            bandeja de tareas y entregables
|   +-- perfil/    (28)            el dossier del usuario
|   +-- firmas/    (8)             firma de PDFs
|   +-- auth/      (8)             login, registro, bootstrap
|   +-- dossier/   (1)             solo el servicio
|   +-- procesos/  (1)
+-- shared/     componentes y utilidades reutilizables
\end{diagrama}

\begin{nota}[Lo que NO existe]
No hay carpeta \rt{src/assets/} (los estaticos viven en \rt{frontend/public/}), ni
\rt{src/views/} ni \rt{src/components/} de nivel raiz: todo esta bajo \rt{modules/} y
\rt{shared/}. Y el \textbf{chat no es un modulo}: vive como widget compartido en
\rt{shared/components/widgets/WorkspaceChatLauncher.vue}.
\end{nota}

\section{El arranque}

\rt{src/main.js} tiene 569 bytes y hace lo minimo: crea la app, registra globalmente un unico
componente (\rt{font-awesome-icon}), importa los dos CSS globales, marca el entorno local en un
atributo del DOM, y monta el router.

\rt{App.vue} solo tiene dos cosas: un \rt{<router-view :key="route.fullPath">} y un
\rt{SessionExpiryModal} global gobernado por \rt{useSessionMonitor}.

Sobre Vite: no hay \rt{server.proxy} --- el proxy es nginx, externo. El alias \rt{@} apunta a
\rt{src}. Y de Tailwind v4 \textbf{no hay \rt{tailwind.config.js} ni \rt{postcss.config.js}}: se
usa el plugin de Vite mas directivas dentro del propio CSS
(\rt{@import "tailwindcss";} en \rt{shared/styles/tailwind.css}).

\section{Estado: no hay Pinia ni Vuex}

Esto sorprende viniendo de cualquier tutorial de Vue. El patron real es de tres capas:

\begin{enumerate}[leftmargin=1.6em]
  \item \textbf{Composables que poseen su propio estado.} Una funcion \rt{useXxx()} declara sus
        \rt{ref} y los devuelve. Ejemplo: \rt{useFlowBuilder.js} tiene seis \rt{ref} propios y
        solo recibe tres dependencias inyectadas. Los comentarios documentan una migracion
        deliberada desde el antipatron ``recibir refs ajenos y escribirlos''.
  \item \textbf{La sesion vive en \rt{localStorage}}: la clave \rt{token} (el JWT) y \rt{user}
        (el usuario serializado como JSON). Las escribe \rt{AuthService.js} y las limpia
        \rt{clearAuthData()} de \rt{core/utils/tokenUtils.js}.
  \item \textbf{Eventos del \rt{window} como bus informal}: \rt{"dossier-updated"},
        \rt{"workspace-chat:context-updated"}, \rt{"workspace-chat:open-process"}.
\end{enumerate}

\begin{aviso}[Consecuencia real que hay que conocer]
Como el usuario esta en \rt{localStorage} y no en un store reactivo, \textbf{los cambios de
sesion no propagan por Vue}. De ahi el truco \rt{<router-view :key="route.fullPath">} en
\rt{App.vue}, que fuerza a remontar el componente entero en cada navegacion.
\end{aviso}

\section{La capa HTTP}

\rt{src/core/services/httpClient.js} es una instancia \emph{propia} de axios (no el singleton
global; el comentario del fichero lo documenta como correccion de un bug de orden de imports)
con \textbf{un solo interceptor de request}, que anade
\rt{Authorization: Bearer <localStorage.token>} si no viene ya puesta.

\begin{aviso}[No hay interceptor de response]
Es decir, \textbf{no hay refresh automatico ante un 401}. El refresh es manual y explicito:
\rt{useSessionMonitor.js} hace polling cada 60 segundos y, cinco minutos antes de que expire el
JWT, abre \rt{SessionExpiryModal.vue}, que llama a \rt{/users/refresh-token} con
\rt{withCredentials: true} (para que viaje la cookie \rt{refreshToken}). Si falla, cierra sesion.
\end{aviso}

Las URLs se construyen en \rt{src/core/config/apiConfig.js}, que expone un objeto
\rt{API\_ROUTES} con unas 110 entradas (unas son cadenas, otras funciones):

\begin{lstlisting}[language=Java]
const API_PORT = import.meta.env.VITE_API_PORT || "3030";
const RAW_API_BASE_URL = String(import.meta.env.VITE_API_BASE_URL || "").trim();
const API_BASE_URL = RAW_API_BASE_URL || `${protocol}//${hostname}:${API_PORT}`;
export const API_PREFIX = `${NORMALIZED_API_BASE_URL}/deasy/v1`;   // -> "/api/deasy/v1"
\end{lstlisting}

Regla documentada en el propio fichero: \emph{``en \rt{frontend/src} no se importa \rt{axios} a
pelo''}. Y no se fija \rt{baseURL} a proposito: todas las llamadas pasan la URL completa
construida desde \rt{apiConfig}.

Hay 17 ficheros de servicio, todos clases con metodos asincronos sobre \rt{httpClient}. El mas
grande es \rt{core/services/ProcessDefinitionPanelService.js} (16,5 KB).

\section{El router y sus guards}

Todo en \textbf{un solo fichero}: \rt{src/core/router/index.js}. Doce rutas de primer nivel mas
ocho hijas bajo \rt{/perfil}. \textbf{Sin lazy-loading}: todas las vistas se importan
estaticamente (el \rt{import()} diferido solo se usa para los dos grafos de admin).

El \rt{router.beforeEach} hace, en este orden:

\begin{enumerate}[leftmargin=1.6em]
  \item Consulta el estado de bootstrap del sistema (?`esta instalado?). Si el modo no es
        \rt{normal}, limpia la sesion y fuerza \rt{/setup}. Si la llamada falla, avisa por
        consola y \textbf{deja pasar} (que el backend este caido no bloquea el enrutado).
  \item Si no hay token valido y la ruta es privada, limpia la sesion y va a \rt{/}.
  \item Si el usuario es administrador y la ruta tiene \rt{meta.blockedForAdmin}, redirige a
        \rt{/admin}.
  \item Comprueba \rt{meta.requiresProcessManagementAccess} y \rt{meta.requiresAdminAccess}.
\end{enumerate}

\begin{concepto}[Por que se usa \rt{meta} y no nombres de ruta]
Porque vue-router \textbf{hereda el \rt{meta} a las rutas hijas}. \rt{/perfil} tiene ocho hijas
(\rt{formacion}, \rt{experiencia}, \rt{referencias}, \rt{capacitacion}, \rt{certificacion},
\rt{investigacion}, \rt{certificados-firma}...). Marcando solo el padre con
\rt{blockedForAdmin: true}, las ocho quedan protegidas automaticamente. Si se comprobara por
nombre de ruta, habria que acordarse de anadir cada hija nueva a una lista.
\end{concepto}

Las rutas publicas son \rt{/}, \rt{/register}, \rt{/recover-password}, \rt{/terminos} y
\rt{/setup}. La ruta \rt{/logout} no tiene componente: solo un \rt{beforeEnter} que llama al
backend, limpia y devuelve \rt{/}.

\begin{nota}[Coste a tener en cuenta]
El guard hace una peticion HTTP (el estado de bootstrap) antes de \emph{cada} navegacion, sin
cache.
\end{nota}

\rt{src/core/utils/accessControl.js} es un \textbf{espejo declarado} de \rt{RbacService.js} del
backend: \rt{hasAnyRole}, \rt{hasPermission} (donde \rt{AdminSistema} pasa todo),
\rt{canAccessResource(resource, action)} con respaldo a \rt{.manage}, y un mapa de unas 50
tablas a su recurso RBAC.

\begin{aviso}[El RBAC del frontend es solo cosmetico]
Sirve para \emph{ocultar botones} y no ensenar opciones que van a fallar. La autorizacion real
la hace \textbf{siempre} el backend. Nunca confies en un control de acceso que corre en el
navegador del usuario: es codigo que el usuario controla.
\end{aviso}

\section{Composables destacados}

El patron comun: una funcion \rt{useXxx(\{ deps \})} que \textbf{declara y devuelve sus propios
\rt{ref}}, con dependencias inyectadas de solo lectura. Nada de \rt{provide}/\rt{inject}, nada
de singletons reactivos.

\rt{useFlowBuilder.js} (modulo \rt{home}) es el constructor del \emph{flujo routed}: quien
elabora (\rt{flowEntrega}) y quien firma (\rt{flowFirma}, con pasos ordenados que llevan
\rt{signers}, \rt{approval\_mode} y \rt{required\_min}). Fue extraido de \rt{HomeView.vue}
durante la fase B del refactor del God Object.

Otros: \rt{useProcessPanels.js}, \rt{useDeliverableView.js} (41 KB, el mayor),
\rt{useDocumentCenter.js}, \rt{useGeneralTask.js}, \rt{useRecipientSearch.js},
\rt{useWorkspaceChrome.js} (estado del \emph{chrome} compartido por las cuatro vistas que montan
\rt{AppWorkspaceShell}), y 24 composables en el modulo admin repartidos en seis subcarpetas
(\rt{data/}, \rt{fk/}, \rt{forms/}, \rt{modals/}, \rt{processes/}, \rt{ui/}).

\section{Tiempo real en el cliente}

Un unico cliente: \rt{src/core/services/realtimeClient.js} (clase \rt{RealtimeClient} exportada
como singleton).

\begin{itemize}[leftmargin=1.4em]
  \item Autentica con \rt{auth: \{ token: localStorage.token \}} --- el mismo JWT que el
        \rt{httpClient}.
  \item Resuelve el destino segun la base: si es absoluta, \rt{io(base, \{path: "/socket.io"\})};
        si es relativa (\rt{/api}, detras de nginx), \rt{io(\{path: "/api/socket.io"\})} sobre
        el mismo origen.
  \item Reconexion infinita con backoff de 1 a 8 segundos; \rt{connect()} es idempotente y
        \textbf{reconecta si el token cambio}.
\end{itemize}

Su unico consumidor es \rt{WorkspaceChatLauncher.vue}, que escucha
\rt{chat.message.created} y \rt{chat.notification.created}.

\begin{nota}[Codigo muerto detectado]
\rt{subscribeProcess} y \rt{unsubscribeProcess} estan definidos en el cliente pero
\textbf{no tienen ningun llamador} en \rt{src/}.
\end{nota}

\section{Vue Flow y dagre: los dos grafos}

Se usan \textbf{solo} en el modulo admin, para dos grafos interactivos:

\begin{itemize}[leftmargin=1.4em]
  \item \textbf{Organigrama de unidades} --- \rt{modules/admin/components/units/UnitGraphView.vue}
        (57 KB), con nodos y aristas propios (\rt{UnitNode.vue}, \rt{UnitEdge.vue}).
  \item \textbf{Mapa de procesos} --- \rt{ProcessGraphView.vue} (55 KB), con \rt{ProcessNode.vue},
        \rt{ProcessConfigNode.vue} y \rt{ProcessTemplateNode.vue}.
\end{itemize}

\rt{dagre} calcula el layout jerarquico de arriba a abajo automaticamente, asi no hay que
posicionar nodos a mano. Ambos se cargan con \rt{defineAsyncComponent(() => import(...))}, de
modo que Vue Flow y dagre \textbf{quedan fuera del bundle inicial} y solo se descargan al abrir
esas dos pantallas.

Se alcanzan por URL: \rt{/admin/academia/unidades/organigrama} y
\rt{/admin/gestiones/procesos/mapa}.

\begin{nota}
El constructor de flujos de firma \textbf{no} usa Vue Flow: es un formulario de listas
(\rt{useFlowBuilder.js}).
\end{nota}

\section{Los ``God Objects'' (deuda tecnica visible)}

\begin{center}
\footnotesize
\begin{tabularx}{\textwidth}{@{}LP{2.6cm}@{}}
\toprule
\textbf{Fichero} & \textbf{Tamano} \\
\midrule
\rt{modules/home/views/HomeView.vue}                       & \textbf{236 KB} (unas 5.600 lineas) \\
\rt{modules/admin/components/tables/AdminTableManager.vue}  & 158 KB \\
\rt{modules/firmas/components/FirmarPdf.vue}                & 120 KB \\
\rt{modules/admin/components/modals/AdminDraftArtifactModal.vue} & 77 KB \\
\rt{modules/admin/components/units/UnitGraphView.vue}       & 58 KB \\
\rt{modules/admin/components/units/ProcessGraphView.vue}    & 56 KB \\
\rt{modules/firmas/components/MultiSignerPanel.vue}         & 52 KB \\
\rt{modules/admin/views/AdminView.vue}                      & 46 KB \\
\rt{modules/home/composables/useDeliverableView.js}         & 41 KB \\
\bottomrule
\end{tabularx}
\end{center}

Estan en el plan de calidad y se estan troceando extrayendo composables --- asi nacio
\rt{useFlowBuilder.js}.

\begin{concepto}[Que es un ``God Object'']
Un objeto (aqui, un componente) que sabe y hace demasiado: acumula responsabilidades de varios
dominios hasta que nadie puede cambiarlo sin miedo. El sintoma es el tamano, pero el problema
real es que \emph{cualquier} cambio lo toca, asi que todos los cambios entran en conflicto y
ningun test lo cubre entero. La cura no es partirlo por lineas, sino \textbf{extraer
responsabilidades completas} (un composable con su estado propio) una a una.
\end{concepto}

% =====================================================================
\chapter{El microservicio signer}
% =====================================================================

Es la pieza mas especializada del sistema: Python 3.14 mas \textbf{pyHanko}, la biblioteca de
firma de PDF.

\section{Arquitectura interna}

\rt{signer/app.py} (1.545 lineas) arranca \textbf{dos cosas a la vez}:

\begin{lstlisting}[language=Python]
rabbit_thread = threading.Thread(target=start_rabbit_worker, daemon=True)
rabbit_thread.start()      # <- el camino REAL de produccion
start_http_server()        # <- solo para depurar (POST /sign, /validate, GET /health)
\end{lstlisting}

El worker de RabbitMQ escucha dos colas (\rt{deasy.sign.request} y
\rt{deasy.sign.validate.request}) con \rt{prefetch\_count=1}: \textbf{un trabajo a la vez}, para
no saturar la memoria con PDFs grandes. Ante cualquier excepcion, registra el error, duerme 5
segundos y reintenta la conexion en un bucle infinito.

Un detalle de higiene: la URL de conexion pasa por \rt{redact\_amqp\_url()} antes de aparecer en
los logs, sustituyendo la contrasena por asteriscos.

\section{El flujo completo de una firma, paso a paso}

\begin{enumerate}[leftmargin=1.8em]
  \item El frontend llama a \rt{/api/deasy/v1/sign/...}
  \item El backend genera un \rt{correlationId}, crea la cola de respuesta
        \rt{deasy.sign.response.<uuid>}, publica el trabajo y entra en polling.
  \item El signer recibe el mensaje y \textbf{valida el payload}: exige \rt{signType},
        \rt{minioPdfPath}, \rt{stampText}, \rt{finalPath}, \rt{minioCertPath} y
        \rt{certPassword}; \rt{signType} debe ser \rt{coordinates} o \rt{token}.
  \item Crea un directorio temporal con permisos \rt{0700} \textbf{fuera de \rt{/tmp}} ---
        porque ahi se desempaqueta el certificado \rt{.p12} descifrado.
  \item Descarga de MinIO el PDF y el certificado.
  \item Abre el \rt{.p12} y \textbf{rechaza el trabajo si el certificado esta expirado o aun no
        es valido}, con mensajes explicitos en espanol.
  \item Si el PDF \textbf{ya tiene firmas embebidas}, salta el paso de limpieza con Ghostscript
        --- aplanarlo destruiria las firmas existentes.
  \item Resuelve donde firmar: coordenadas explicitas, o busca un \textbf{token de texto} dentro
        del PDF con \rt{pdfplumber} (\rt{find\_marker.py} recorre las palabras de cada pagina y
        devuelve sus posiciones).
  \item \textbf{Llama a Node} para generar la imagen del sello (ver seccion siguiente).
  \item Por cada coordenada, pyHanko crea un campo \rt{sig\_<uuid>} e inserta la firma: PAdES,
        SHA-256, con sellado de tiempo opcional. Con varias firmas, \textbf{cada iteracion parte
        del PDF firmado anterior} (firmas encadenadas).
  \item Valida la firma resultante contra el contexto de confianza construido desde
        \rt{signer/trust/}, y aborta salvo tolerancia explicita en el payload.
  \item Sube el PDF firmado a MinIO, reemplazando \rt{minioPdfPath}, y tambien a \rt{finalPath}
        si es distinto.
  \item Publica el resultado en la cola de respuesta; el backend, que estaba haciendo polling,
        lo recoge.
  \item Al salir del bloque \rt{with}, \textbf{todo el material temporal se borra}, incluido el
        \rt{.p12}.
\end{enumerate}

\section{\rt{sigmaker/}: por que hay Node dentro de un servicio Python}

Porque la imagen del sello (QR + nombre del firmante + logo de la PUCE + fecha) se genera con
las bibliotecas \rt{canvas} y \rt{qrcode} de Node, mejores para eso que las alternativas de
Python. Python lo invoca como subproceso:

\begin{lstlisting}[language=Python]
subprocess.run(["node", SIGMAKER_DIR + "/cli.mjs",
                stamp_png, stamp_text, final_path, logo_path], check=True)
\end{lstlisting}

El QR codifica \textbf{la URL publica del documento firmado}, para que quien reciba el PDF
impreso pueda verificarlo escaneandolo. La fecha se calcula en la zona horaria
\rt{America/Guayaquil} y el PNG se escribe con fondo transparente.

Por eso el \rt{docker/signer/Dockerfile} instala \rt{nodejs} y \rt{npm} dentro de una imagen
Python, y hace un \rt{npm install} separado en \rt{/opt/sigmaker}.

\section{\rt{trust/}: la cadena de confianza}

\begin{diagrama}
trust/roots/   <- CAs raiz acreditadas en Ecuador
                  firmasegura-root-ca-1.pem
                  icert-ec-root.pem
                  security-data-root-ca-2.pem

trust/extra/   <- sus subordinadas (se cargan como other_certs, NO como anclas)
                  firmasegura-subca-1.pem
                  icert-ec-subordinate.pem
                  security-data-subca-2.pem
\end{diagrama}

El contexto de validacion se construye con \rt{allow\_fetching=True} y
\rt{revocation\_mode="soft-fail"}.

\begin{concepto}[Por que ``raiz'' y ``subordinada'' van en sitios distintos]
Una \emph{ancla de confianza} (trust anchor) es un certificado que decides creer sin
demostracion: es el final de la cadena. Las subordinadas, en cambio, son eslabones intermedios:
hay que \emph{poder verificarlas} contra una raiz, no confiar en ellas por decreto. Si metieras
una subordinada como ancla, aceptarias como valida una cadena que en realidad no llega a
ninguna raiz acreditada.
\end{concepto}

\section{Los OIDs ecuatorianos: 500 lineas de \rt{app.py}}

Casi un tercio del fichero esta dedicado a \textbf{extraer la cedula y el nombre del firmante}
de los certificados, porque \textbf{cada emisor ecuatoriano usa OIDs propietarios distintos}:

\begin{center}
\footnotesize
\begin{tabularx}{\textwidth}{@{}P{4.0cm}L@{}}
\toprule
\textbf{Emisor} & \textbf{Prefijo de OID} \\
\midrule
Security Data & \rt{1.3.6.1.4.1.37746.3.*} \\
ICERT         & \rt{1.3.6.1.4.1.43745.1.3.*} \\
Firmasegura   & \rt{1.3.6.1.4.1.61305.3.*} \\
\bottomrule
\end{tabularx}
\end{center}

No hay estandar comun: hay que conocer los tres. Ademas se leen por \textbf{dos vias distintas}
(\rt{cryptography} y \rt{asn1crypto} sobre el TBS crudo) porque ninguna cubre todos los casos.

\section{Pruebas del signer}

Todas en \rt{signer/tests/}, con \rt{unittest} de la biblioteca estandar (no hay pytest en la
imagen). \textbf{266 metodos \rt{test\_}} repartidos en nueve ficheros:

\begin{center}
\footnotesize
\begin{tabularx}{\textwidth}{@{}P{5.4cm}P{1.4cm}L@{}}
\toprule
\textbf{Fichero} & \textbf{Casos} & \textbf{Que fija} \\
\midrule
\rt{test\_certificate\_parsing.py}    & 72 & DN, extensiones, cedula: quien firmo \\
\rt{test\_signing\_flow.py}           & 36 & Caja del sello, offset del modo token, firmas encadenadas \\
\rt{test\_payload\_validation.py}     & 34 & Orden de guards y mensajes de error \\
\rt{test\_pdf\_metadata.py}           & 33 & Lectura del diccionario \rt{/Sig}, fechas \\
\rt{test\_transport.py}               & 24 & Servidor HTTP y publicacion por RabbitMQ \\
\rt{test\_validation\_report.py}      & 21 & Forma exacta del informe que consume el frontend \\
\rt{test\_plumbing.py}                & 20 & MinIO, material de confianza, Ghostscript, temporales \\
\rt{test\_certificate\_extensions.py} & 17 & Las dos fuentes de extensiones, con certificados reales \\
\rt{test\_certificate\_loading.py}    & 9  & Carga del PKCS\#12 y guards de vigencia \\
\bottomrule
\end{tabularx}
\end{center}

Como apoyo hay \rt{tests/dobles.py} (objetos ``pato'' que solo exponen lo que \rt{app.py}
consulta) y \rt{tests/herramientas.py}, que \textbf{no son pruebas}: miden complejidad
ciclomatica y anidamiento via AST.

\begin{nota}[Los conteos no cuadran entre si]
\rt{signer/README.md} dice 224 casos, \rt{requirements-dev.txt} y \rt{sonar-project.properties}
dicen 229, y el conteo real hoy es 266. Los documentos van por detras del codigo.
\end{nota}

% =====================================================================
\chapter{Infraestructura y despliegue}
% =====================================================================

\section{Docker: base mas overlays}

En vez de un \rt{docker-compose.yml} por entorno (duplicando todo), hay \textbf{un fichero base
y un overlay} que lo modifica:

\begin{diagrama}
dev     = compose.base.yml + compose.proxy.yml + compose.dev.yml
qa      = compose.base.yml + compose.qa.yml
prod    = compose.base.yml + compose.prod.yml
ingress = compose.ingress.yml              (stack independiente)
ingress-bootstrap = compose.ingress.bootstrap.yml   (solo HTTP, para el desafio ACME)
\end{diagrama}

\begin{concepto}[Que es un ``overlay'' de compose]
Docker Compose permite pasar varios \rt{-f}. Los ficheros se \textbf{fusionan} en orden: lo que
define el segundo pisa o completa lo del primero. Asi, el base declara los servicios y sus
variables comunes, y cada overlay solo dice lo que cambia (puertos, si se construye la imagen o
se descarga, montajes de codigo, politicas de reinicio). Evita mantener cuatro ficheros casi
identicos que se desincronizan.
\end{concepto}

Consecuencia concreta: \rt{qa} y \rt{prod} \textbf{no cargan} \rt{compose.proxy.yml}, asi que
no tienen \rt{nginx-proxy}; su entrada publica es el stack \rt{ingress}, separado y compartido.

\subsection{Diferencias por entorno}

\begin{center}
\footnotesize
\begin{tabularx}{\textwidth}{@{}P{3.2cm}LL@{}}
\toprule
 & \textbf{dev} & \textbf{qa / prod} \\
\midrule
Imagenes de la app  & se \textbf{construyen} localmente (\rt{Dockerfile.dev}) & se \textbf{descargan} de GHCR con \rt{pull\_policy: always} \\
Codigo              & \textbf{bind mount} (editas y se recarga) & dentro de la imagen \\
Proxy               & \rt{nginx-proxy} en el propio stack & un stack \rt{ingress} separado \\
Puertos del proxy   & 8088 / 8443 & el ingress publica 80 / 443 \\
PostgreSQL          & 5432 & 15432 (qa) / 25432 (prod) \\
MinIO               & 9000 / 9001 & 9100 / 9101 y 9200 / 9201 \\
Buckets             & \rt{deasy-*} & \rt{deasy-qa-*} y \rt{deasy-prod-*} \\
PAdES-LTA           & desactivado & \textbf{activado en prod} \\
Endurecimiento      & --- & \rt{no-new-privileges}, \rt{read\_only} en frontend, \rt{tmpfs}, healthchecks propios \\
\bottomrule
\end{tabularx}
\end{center}

\begin{aviso}[RabbitMQ solo se publica en dev, y atado a \rt{127.0.0.1}]
El comentario del compose explica el motivo: la interfaz de gestion de RabbitMQ
\textbf{muestra el cuerpo de los mensajes}, y ahi viaja \rt{certPassword}, la contrasena del
certificado de firma. En qa y prod los puertos no se publican en absoluto.
\end{aviso}

\subsection{Los servicios del fichero base}

\begin{center}
\footnotesize
\begin{tabularx}{\textwidth}{@{}P{2.9cm}P{4.4cm}L@{}}
\toprule
\textbf{Servicio} & \textbf{Imagen} & \textbf{Notas} \\
\midrule
\rt{postgres}        & \rt{postgres:17} & healthcheck \rt{pg\_isready}; unico motor de datos \\
\rt{rabbitmq}        & \rt{rabbitmq:3-management} & healthcheck \rt{rabbitmq-diagnostics ping} \\
\rt{minio}           & \rt{minio/minio:...} & \rt{server /data --console-address ":9001"} \\
\rt{minio-bootstrap} & \rt{minio/mc:...} & perfil \rt{storage-init}; crea los buckets y hace \rt{mc mirror} de las semillas \\
\rt{backend}         & sin imagen en base & la define cada overlay \\
\rt{frontend}        & sin imagen en base & idem \\
\rt{signer}          & sin imagen en base & \rt{SIGNER\_PORT=4000}, \rt{SIGMAKER\_DIR=/opt/sigmaker} \\
\rt{analytics}       & sin imagen en base & perfil \rt{workers}; hoy es un marcador de posicion \\
\bottomrule
\end{tabularx}
\end{center}

\section{\rt{scripts/docker-env.sh}: la interfaz comun}

Es \textbf{la unica forma correcta} de tocar el stack. Todo lo que va despues del entorno se
reenvia literalmente a \rt{docker compose}:

\begin{lstlisting}[language=bash]
bash scripts/docker-env.sh dev up -d --build
bash scripts/docker-env.sh dev logs -f backend
bash scripts/docker-env.sh dev ps
bash scripts/docker-env.sh dev exec -T backend npm run test:unit
bash scripts/docker-env.sh dev down
\end{lstlisting}

Resuelve el fichero \rt{.env} correcto (con \rt{DEASY\_ENV\_FILE} como gancho para apuntar al
\rt{.runtime} en despliegues), valida que existan todos los ficheros antes de invocar nada,
exporta \rt{DEASY\_CONTAINER\_ENV\_FILE} y hasta convierte rutas si detecta \rt{cygpath} en
Windows.

\begin{aviso}[Regla del proyecto]
\emph{``Builds y tests deben correr dentro de los contenedores via \rt{scripts/docker-env.sh},
no con npm o npx en el host''}. Y al \textbf{anadir dependencias al frontend} hay que instalar
\textbf{dentro del contenedor}: el volumen de \rt{node\_modules} sombrea el de la imagen, asi
que un \rt{pnpm install} en el host \textbf{no se ve} dentro.
\end{aviso}

\subsection{Los demas scripts operativos}

\begin{center}
\footnotesize
\begin{tabularx}{\textwidth}{@{}P{4.4cm}L@{}}
\toprule
\textbf{Script} & \textbf{Que hace} \\
\midrule
\rt{apply-env.sh}       & Motor comun de despliegue. Aplica el tag de imagen, crea la red externa si falta, limpia contenedores antiguos que retengan los puertos 80/443, y levanta el stack. Soporta \rt{DEASY\_DRY\_RUN=1}. \\
\rt{deploy-env.sh}      & Envoltorio fino que llama a \rt{apply-env.sh} en modo \rt{gh-actions}. Es lo que ejecuta el workflow por SSH. \\
\rt{server-pull-deploy.sh} & Modo \emph{server-pull}: exige rama correcta y arbol limpio, hace \rt{git pull --ff-only} y despliega. \\
\rt{\_backend\_db\_exec.sh} & Libreria compartida. Para produccion exige un fichero de aprobacion \textbf{dentro del repo e ignorado por git}. \\
\rt{reset-db.sh}        & Resetea solo la base de datos dentro del contenedor. \\
\rt{reset-system.sh}    & Reset completo a instalacion virgen: base + almacenamiento, recicla backend y signer, y \textbf{recarga nginx} (porque el proxy de dev cachea la IP del backend). \\
\rt{rabbitmq-migrar-usuario.sh} & Crea el usuario real de RabbitMQ y \textbf{despues borra \rt{guest}}. Idempotente. \\
\rt{bootstrap-ingress-cert.sh} & Emite el certificado publico con certbot (ver mas abajo). \\
\rt{signer-coverage.sh} & Genera el informe de cobertura de Python para Sonar, reescribiendo las rutas para que las resuelva. \\
\bottomrule
\end{tabularx}
\end{center}

\section{nginx: dos configuraciones distintas}

\subsection{Mecanica de plantillas}

\textbf{No existe un directorio \rt{nginx/conf.d/} en el repositorio.} Los tres directorios
\rt{*-conf.d/} se montan como \rt{/etc/nginx/templates}, y el entrypoint oficial de la imagen
\rt{nginx} aplica \rt{envsubst} a cada \rt{*.template} generando los \rt{.conf} finales. De ahi
que las plantillas usen variables \rt{\$\{VAR\}} literales.

\subsection{Desarrollo}

\rt{nginx/app-conf.d/default.conf.template}: un \rt{server\_name \_} que no discrimina.

\begin{lstlisting}[language=bash]
location = /api   { return 301 /api/; }
location /api/    { proxy_pass http://${BACKEND_UPSTREAM}/; }
location /        { proxy_pass http://${FRONTEND_UPSTREAM}; }
\end{lstlisting}

\begin{concepto}[La barra final del \rt{proxy\_pass}]
Es la diferencia entre que funcione y que no. Con \rt{proxy\_pass http://backend:3030/;} (con
barra), nginx \textbf{sustituye} la parte que casó (\rt{/api/}) por la barra, asi que
\rt{/api/deasy/v1/users} llega al backend como \rt{/deasy/v1/users}. \emph{Sin} la barra final,
la ruta se pasaria integra y el backend recibiria \rt{/api/deasy/v1/users}, que no tiene
ninguna ruta registrada: 404.
\end{concepto}

El puerto 80 sirve solo tres cosas: \rt{/healthz}, el desafio ACME
(\rt{/.well-known/acme-challenge/}) y una redireccion 301 a HTTPS.

\subsection{Produccion y QA}

\rt{nginx/ingress-conf.d/default.conf.template} es un \textbf{ingress compartido} que reparte
\textbf{por cabecera \rt{Host}}: \rt{fresvel.com} y \rt{www.fresvel.com} van a prod,
\rt{qa.fresvel.com} va a qa. Un cuarto bloque \rt{default\_server} devuelve \textbf{404} a
cualquier Host no reconocido. Aqui la reescritura se hace con
\rt{rewrite \^{}/api/(.*)\$ /\$1 break;}.

Detalle importante: los upstreams se asignan a \textbf{variables} junto con
\rt{resolver 127.0.0.11 ipv6=off valid=10s}. Eso fuerza a nginx a \textbf{re-resolver el DNS de
Docker cada 10 segundos} en vez de cachear la IP al arrancar. Es justo el problema que sufre el
proxy de dev (resolucion estatica), y por eso \rt{reset-system.sh} tiene que recargarlo tras
reciclar el backend.

\subsection{Rutas efectivas}

nginx \textbf{solo conoce dos prefijos}: \rt{/api/} y \rt{/}. Todo lo demas son rutas internas
de Express:

\begin{diagrama}
https://host/api/deasy/v1/users   ->  nginx quita /api  ->  backend /deasy/v1/users
https://host/api/deasy/docs       ->  nginx quita /api  ->  backend /deasy/docs   (Swagger UI)
https://host/api/socket.io        ->  nginx quita /api  ->  backend /socket.io    (WebSocket)
\end{diagrama}

El upgrade a WebSocket funciona gracias a las cabeceras \rt{Upgrade} y \rt{Connection} y al
\rt{map \$http\_upgrade \$connection\_upgrade} definido en \rt{nginx/nginx.conf}.

\subsection{Certificados y ACME}

\begin{itemize}[leftmargin=1.4em]
  \item \rt{nginx/certs/dev/} y \rt{nginx/certs/qa/}: \textbf{autofirmados versionados} en git
        (generables con \rt{nginx/scripts/generate-self-signed.sh}).
  \item \rt{nginx/certs/public/}: \textbf{no versionado}; ahi deja
        \rt{scripts/bootstrap-ingress-cert.sh} el certificado de Let's Encrypt.
  \item \rt{nginx/acme/}: el \emph{webroot} del desafio HTTP-01; montado de solo lectura en
        proxy e ingress, y de \textbf{lectura-escritura} en el bootstrap, para que certbot pueda
        escribir el fichero de reto.
\end{itemize}

\begin{concepto}[Por que existe un ``ingress-bootstrap'']
Let's Encrypt verifica que controlas el dominio pidiendote colocar un fichero en
\rt{http://tudominio/.well-known/acme-challenge/xxx}. Pero el ingress normal solo escucha en
HTTPS con un certificado... que aun no tienes. Es el problema del huevo y la gallina. La
solucion: un nginx minimo, \textbf{solo HTTP}, que sirve el reto; se emite el certificado, se
copia, se baja el bootstrap y se levanta el ingress definitivo.
\end{concepto}

\section{CI/CD}

\subsection{\rt{cd-multienv.yml}: ocho jobs}

\begin{diagrama}
prepare (rama -> entorno)
   |-- frontend-lint        (pnpm run lint)
   |-- backend-checks       (npm run check:imports + npm run test:unit)
   +-- compose-validate     (docker-env.sh <env> config, para los 4 entornos)
              |
              v
        publish-images      (matriz de 4 imagenes -> GHCR)
              |             tags: :dev|:qa|:prod  y  :sha-XXXXXXX
              v
        deploy-app-push / deploy-app-manual / deploy-ingress
              (rsync de docker/ y nginx/ + SSH, solo si DEPLOY_DELIVERY_MODE == 'gh-actions')
\end{diagrama}

Mapeo de ramas: \rt{develop} $\rightarrow$ dev (publica imagenes, \textbf{no despliega});
\rt{qa} $\rightarrow$ qa; \rt{main} $\rightarrow$ prod. El trabajo se hace en \rt{develop};
\rt{main} es la rama de produccion.

El \rt{rsync} solo sube \rt{docker/}, \rt{nginx/} (excluyendo \rt{certs/}) y cuatro scripts. El
fichero de variables reales llega como secreto de GitHub y se vuelca en
\rt{docker/.env.<env>.runtime} en el servidor.

\subsection{Dos modos de entrega}

\begin{itemize}[leftmargin=1.4em]
  \item \textbf{Push}: GitHub Actions entra por SSH y despliega.
  \item \textbf{Pull}: unidades systemd en el servidor
        (\rt{deploy/systemd/deasy-server-pull@.timer}) que cada \textbf{15 minutos} hacen
        \rt{git pull} y redespliegan (\rt{OnBootSec=5m}, \rt{OnUnitActiveSec=15m},
        \rt{Persistent=true}).
\end{itemize}

El segundo modo existe para servidores \textbf{sin IP publica estable}, donde GitHub no puede
iniciar la conexion.

\subsection{\rt{sonar.yml}: solo manual, y en verde si falta configuracion}

Se dispara solo con \rt{workflow\_dispatch}. Un primer paso \emph{guard} comprueba si existen
los secretos \rt{SONAR\_HOST\_URL} y \rt{SONAR\_TOKEN}; si faltan, escribe una explicacion en el
resumen del workflow y \textbf{termina en verde}.

\begin{nota}[La razon documentada]
El SonarQube del proyecto esta autoalojado en \rt{localhost:9002} y \textbf{un runner de GitHub
no puede alcanzarlo}. Poner el analisis en el workflow de push lo dejaria en rojo para siempre,
lo que entrena a todo el mundo a ignorar los fallos de CI. Sonar en CI esta pendiente de una
decision de infraestructura (publicarlo con TLS o migrar a SonarCloud), no de codigo.
\end{nota}

\section{SonarQube: como esta montado}

\begin{center}
\footnotesize
\begin{tabularx}{\textwidth}{@{}P{2.6cm}L@{}}
\toprule
\textbf{Pieza} & \textbf{Donde} \\
\midrule
Stack     & \rt{scripts/sonar/compose.yml} --- SonarQube \emph{community} + PostgreSQL 16, puerto \textbf{9002} \\
Lanzador  & \rt{scripts/sonar/scan.sh} --- \rt{sonar-scanner-cli} dockerizado \\
Config    & \rt{sonar-project.properties} --- \rt{projectKey=deasy}; fuentes \rt{backend}, \rt{frontend/src}, \rt{signer}, \rt{scripts}, con los tests \textbf{declarados aparte} \\
CI        & \rt{.github/workflows/sonar.yml} \\
\bottomrule
\end{tabularx}
\end{center}

\begin{aviso}[Cuatro cosas que cuestan una medicion entera]
\begin{enumerate}[leftmargin=1.4em,nosep]
  \item \textbf{Regenera los DOS informes de cobertura antes de escanear.} Si no, Sonar lee los
        de la corrida anterior \textbf{sin quejarse}. Y si las rutas \rt{SF:} del lcov no son
        relativas a la raiz del repo, los descarta \textbf{en silencio} y la cobertura vuelve a
        cero.
  \item \textbf{Al consultar la API, filtra con \rt{resolved=false}.} Por defecto incluye las
        cerradas y las \emph{won't fix}, y los conteos salen inflados.
  \item \textbf{El escaneo se procesa despues de subirse.} Consultar metricas justo al terminar
        devuelve las viejas: espera a que la tarea diga \rt{SUCCESS}.
  \item \textbf{No toques \rt{sonar.projectVersion}.} Mueve el periodo de New Code y tira la
        serie historica, que es el unico termometro fiable que hay.
\end{enumerate}
\end{aviso}

\begin{nota}[Marcar no es arreglar]
Los falsos positivos se marcan en Sonar con justificacion; hay varios que \textbf{no} hay que
``corregir''. Sonar rastrea la incidencia por el \textbf{hash de la linea}: un renombrado puro
conserva la marca, pero reescribir la linea la pierde.
\end{nota}

% =====================================================================
\chapter{Testing: dos niveles que no se mezclan}
% =====================================================================

Esta es una regla dura del proyecto, y la ubicacion de los ficheros \textbf{no es estetica}:
\rt{sonar-project.properties} distingue tests de codigo de produccion \textbf{por patron}, y un
fichero mal colocado o mal nombrado se analiza \textbf{como codigo de produccion}. Asi es como
las contrasenas de fixture acabaron contando como vulnerabilidades del sistema.

\begin{center}
\footnotesize
\begin{tabularx}{\textwidth}{@{}P{2.6cm}P{4.6cm}P{3.0cm}L@{}}
\toprule
\textbf{Nivel} & \textbf{Donde} & \textbf{Nombre} & \textbf{Que prueba} \\
\midrule
\textbf{Unitario} &
  \textbf{junto al modulo} que prueba &
  \rt{<modulo>.test.js} (\rt{.test.mjs} si el modulo es ESM \rt{.mjs}) &
  Una unidad, sin red ni base de datos \\
\textbf{Caracterizacion} &
  \rt{backend/tests/characterization/flows/} &
  \rt{<flujo>.test.mjs} &
  Contrato HTTP extremo a extremo contra \emph{goldens} \\
\bottomrule
\end{tabularx}
\end{center}

\section{Reglas de ubicacion}

\begin{itemize}[leftmargin=1.4em]
  \item \rt{*.test.js} y \rt{*.test.mjs} son los \textbf{unicos sufijos validos}. \textbf{No uses
        \rt{*.spec.js}}: no hay ni uno en el repositorio, y anadirlo obliga a mantener un patron
        muerto en la configuracion de Sonar y en los globs de \rt{test:unit}.
  \item Un test unitario nuevo tiene que \textbf{caer dentro de los globs de
        \rt{backend/package.json} $\rightarrow$ \rt{test:unit}} (\rt{config/}, \rt{database/**},
        \rt{services/**}, \rt{utils/**}, \rt{middlewares/**}, \rt{controllers/**},
        \rt{errors/**}). Si lo pones
        fuera, \textbf{no lo ejecuta nadie} y no te vas a enterar.
  \item Si el modulo vive en otra carpeta, \textbf{amplia el glob en el mismo commit y en los DOS
        sitios}: \rt{test:unit} y \rt{test:unit:coverage} llevan la misma lista duplicada (el
        segundo con prefijo \rt{backend/}, porque corre desde la raiz del repo para que las rutas
        del lcov le valgan a Sonar). Si solo tocas uno, el test corre pero \textbf{no cuenta para
        la cobertura}.
  \item No metas tests unitarios en \rt{backend/tests/}: esa carpeta es del harness de
        caracterizacion.
  \item Los tests del frontend van junto al componente o composable (vitest los descubre solo).
\end{itemize}

\begin{aviso}[Una suite que no arranca cuenta como fallo, no como ``0 tests'']
Vitest marca \emph{Failed Suite} con 0 casos cuando el error ocurre \textbf{al importar} ---
tipicamente un \rt{vi.mock} incompleto de un modulo que ahora tira de otro. \textbf{Mira la
linea \rt{Test Files}}, no solo la de \rt{Tests}.
\end{aviso}

\section{Los tests unitarios}

\textbf{Backend}: 32 ficheros, \textbf{523 casos}, con \rt{node --test} --- el runner nativo de
Node, sin Jest ni Vitest. Los mas nutridos:

\begin{center}
\footnotesize
\begin{tabularx}{\textwidth}{@{}LP{1.6cm}@{}}
\toprule
\textbf{Fichero} & \textbf{Casos} \\
\midrule
\rt{config/postgres.test.js}                                & 67 \\
\rt{config/postgres.dialect.test.js}                        & 38 \\
\rt{services/admin/templates/workflows.test.js}             & 37 \\
\rt{services/admin/crud/validation.test.js}                 & 33 \\
\rt{services/admin/templates/flowRows.test.js}              & 31 \\
\rt{services/documents/FillRequestWorkflowService.test.js}  & 30 \\
\rt{database/postgres\_schema.test.js}                      & 25 \\
\rt{services/sign/PdfSigningService.test.js}                & 24 \\
\rt{services/sign/BatchSigningService.test.js}              & 21 \\
\rt{controllers/users/user\_controler.primitives.test.js}   & 21 \\
\bottomrule
\end{tabularx}
\end{center}

\textbf{Frontend}: 18 ficheros, \textbf{304 casos} con Vitest y \rt{@vue/test-utils}. La
configuracion vive en el bloque \rt{test:} de \rt{vite.config.js} (no hay \rt{vitest.config.js}).
El entorno por defecto es \rt{node}; los once tests que montan componentes ponen la pragma
\rt{// @vitest-environment jsdom} en la primera linea. \textbf{No hay pruebas E2E} (ni Playwright
ni Cypress).

\section{Los characterization tests (golden master)}

\begin{concepto}[Que es un golden master]
En vez de escribir ``espero que devuelva \rt{\{x: 1\}}'', \textbf{grabas} lo que devuelve hoy y
lo guardas en un JSON. En cada ejecucion comparas contra el JSON grabado. Si cambia, el test
falla y \textbf{el diff de git es la evidencia} de lo que cambiaste.

Sirve para \textbf{blindar refactors} en codigo que no entiendes del todo: no documentas lo que
\emph{deberia} hacer, sino lo que \emph{hace}, rarezas incluidas. Es la red de seguridad ideal
cuando heredas un sistema.
\end{concepto}

En Deasy son \textbf{19 suites, 204 tests y 281 claves de snapshot}, sin dependencias externas:
\rt{fetch} nativo mas \rt{node:test} y \rt{node:assert/strict}.

\subsection{Anatomia del harness}

\begin{center}
\footnotesize
\begin{tabularx}{\textwidth}{@{}P{3.0cm}L@{}}
\toprule
\textbf{Pieza} & \textbf{Que hace} \\
\midrule
\rt{config.mjs} &
  Todo parametrizable por entorno: \rt{API\_PREFIX}, \rt{BASE\_URL}, timeouts,
  \rt{SNAPSHOT\_MODE} (\rt{compare} o \rt{update}), los tres usuarios de referencia y la
  \rt{FIXTURE} con los IDs esperados. \\
\rt{lib/http.mjs} &
  \rt{request/get/post/put/patch/del}; devuelve siempre \rt{\{status, ok, headers, body\}} y
  \textbf{nunca lanza por codigo de estado}. \\
\rt{lib/auth.mjs} &
  \rt{login(user)} y \rt{tokenFor(userKey)} con cache por ejecucion. \\
\rt{lib/snapshot.mjs} &
  El motor: \rt{matchSnapshot(suite, key, actual)}. En modo \rt{update} escribe y pasa; en
  \rt{compare} falla, \textbf{y tambien falla si la clave no existe}. \\
\rt{lib/normalize.mjs} &
  Enmascara lo volatil (timestamps, \rt{iat}, \rt{exp}, tokens) con \rt{"<normalized>"}. Los IDs
  \textbf{se conservan} a proposito. \\
\rt{lib/readiness.mjs} &
  Polling cada 2 segundos hasta que el sistema responde. \\
\rt{lib/db.mjs} &
  \textbf{Excepcion declarada}: pool de PostgreSQL directo \textbf{solo para limpieza}. La regla
  escrita es \emph{``se usa para LIMPIAR, nunca para ASERTAR''}. \\
\rt{setup/bootstrap\_system.mjs} &
  Construye la fixture con el \textbf{bootstrap real por HTTP}, no con un seed SQL paralelo (que
  era la fuente de verdad anterior y derivaba). Idempotente, y falla si los IDs divergen. \\
\rt{setup/seed\_execution.mjs} &
  Siembra datos de ejecucion via API. \\
\rt{\_\_snapshots\_\_/} &
  19 JSON, uno por suite, con claves ordenadas alfabeticamente e indentacion 2, para que el
  \textbf{diff de git sea legible}. \\
\bottomrule
\end{tabularx}
\end{center}

\subsection{El ciclo de trabajo}

\begin{lstlisting}[language=bash]
test:char:fixture  = reset.mjs db --yes && test:char:bootstrap && test:char:seed
test:char:run      = test:char:fixture && test:char
test:char:capture  = test:char:fixture && SNAPSHOT_MODE=update npm run test:char
test:char          = node --test --test-concurrency=1 tests/characterization/flows/*.test.mjs
\end{lstlisting}

Los prefijos de \rt{z} en los nombres de fichero \textbf{fuerzan el orden alfabetico} de
ejecucion (el runner corre con \rt{--test-concurrency=1}), poniendo al final las suites que
\emph{escriben} datos. Hoy son \textbf{nueve}, y la escalera llega hasta siete \rt{z}:
\rt{zz\_default\_process\_routed}, \rt{zz\_task\_generation}, \rt{zz\_template\_lifecycle},
\rt{zzz\_artifact\_draft}, \rt{zzzz\_sign\_batch}, \rt{zzzz\_sign\_workflow},
\rt{zzzzz\_task\_item\_relay}, \rt{zzzzzz\_flow\_steps\_db} y
\rt{zzzzzzz\_schema\_flow\_reread}.

\begin{aviso}[\rt{test:char:run} RESETEA la base de dev]
Hace reset, bootstrap y seed. Es lo normal para caracterizacion, pero \textbf{no lo lances si
tienes datos que quieras conservar}. Para actualizar los goldens: \rt{test:char:capture}.
\end{aviso}

\begin{aviso}[La regla de oro]
\textbf{Un refactor no cambia un golden.} Si un snapshot se mueve durante un refactor puro, o
rompiste algo o el test estaba mal. En un \emph{fix}, el diff del golden \textbf{es} la prueba
del arreglo.
\end{aviso}

\section{Comandos de validacion por modulo}

\begin{lstlisting}[language=bash]
# Frontend: lint + tests unitarios, y ademas verificar en el navegador
bash scripts/docker-env.sh dev exec -T frontend pnpm run lint
bash scripts/docker-env.sh dev exec -T frontend pnpm run test:unit
bash scripts/docker-env.sh dev exec -T frontend pnpm run test:unit:coverage
bash scripts/docker-env.sh dev exec -T frontend pnpm run build

# Backend: no tiene lint, pero SI tiene tests
bash scripts/docker-env.sh dev exec -T backend npm run test:unit
bash scripts/docker-env.sh dev exec -T backend npm run test:char:run
bash scripts/docker-env.sh dev exec -T backend npm run check:imports   # OBLIGATORIO tras mover codigo
bash scripts/docker-env.sh dev exec -T backend npm run test:unit:coverage
\end{lstlisting}

\section{Las reglas al mover codigo}

Del \rt{CLAUDE.md}, porque cada una nacio de un incidente real:

\begin{enumerate}[leftmargin=1.6em]
  \item \textbf{Refactor = mover codigo, NO reescribir comportamiento.} Si cambias que hace algo,
        no es un refactor.
  \item \textbf{\rt{node --check} valida SINTAXIS, no imports.} Por eso \rt{check:imports} es
        obligatorio, y comprobar que el backend arranca \textbf{no} lo sustituye.
  \item \textbf{No injertes casos especiales en el camino generico.}
  \item \textbf{Extrae por script, no a mano}, y verifica \rt{count == 1} antes de borrar cada
        bloque.
  \item \textbf{El SQL no lo valida nadie} hasta que se ejecuta esa rama. Pruebalo con
        \rt{PREPARE} en psql.
\end{enumerate}

% =====================================================================
\chapter{Cosas que te van a confundir (y no son bugs)}
% =====================================================================

Esta lista ahorra horas. Son incoherencias reales del repositorio, verificadas.

\begin{enumerate}[leftmargin=1.8em,itemsep=6pt]
  \item \textbf{\rt{backend/README.md} esta obsoleto.} Habla de MongoDB (\rt{URI\_MONGO}),
        MariaDB y puerto 3000. El codigo usa PostgreSQL y el puerto 3030.
        \textbf{El \rt{CLAUDE.md} es la fuente fiable}, no los README.

  \item \textbf{El \rt{README.md} raiz enlaza documentos que ya no existen}
        (\rt{docs/02-dominio-datos/modelo-datos.md}, \rt{MER\_SQL.sql},
        \rt{docs/01-arquitectura/}). Estan en \rt{docs/docs-md-antiguos/}. Hoy \rt{docs/02-dominio-datos/}
        tiene \rt{consolidado.dbml} (53 KB, 67 tablas y 139 relaciones), la carpeta
        \rt{dominios/} con los ocho DBML por dominio, y \rt{anotaciones.json}. Los dos primeros
        son \textbf{artefactos generados} por \rt{scripts/docs/gen-dbml.sh}; el unico que se
        escribe a mano es \rt{anotaciones.json}.

  \item \textbf{\rt{amqplib} esta en \rt{package.json} pero no se importa en ningun sitio.} Es
        una dependencia muerta: la integracion real con RabbitMQ es por su API HTTP de gestion.

  \item \textbf{\rt{sqlTables.js} declara campos de \rt{template\_artifacts} que ya no son
        columnas de esa tabla} (\rt{template\_code}, \rt{display\_name}, \rt{description},
        \rt{template\_scope}, \rt{template\_seed\_id}, \rt{owner\_person\_id}). \textbf{No es un
        bug}: \rt{SqlAdminService} los rutea por JOIN a \rt{deliverables}, exponiendolos con los
        mismos nombres (\rt{d.code AS template\_code}). Es una fachada deliberada.

  \item \textbf{``Entregable'' significa dos cosas} segun el contexto: \rt{task\_items} (la
        instancia) y \rt{deliverables} (la identidad de la plantilla).

  \item \textbf{Las fases del plan de calidad se llaman A, B, C... por el orden en que se
        descubrieron}, no por prioridad ni por tema.

  \item \textbf{\rt{test:char:run} resetea la base de dev.}

  \item \textbf{\rt{rounded-lg} vale 16px} en este frontend, no lo que dice la documentacion de
        Tailwind: \rt{theme.css} pisa el espacio de nombres de radios de Tailwind. La escala esta
        efectivamente invertida.

  \item \textbf{La contrasena del gestor NO es \rt{Demo1234!}}, es \rt{Gestor1234!}.

  \item \textbf{No existe \rt{artifact\_stage}} pese a estar descrito en la documentacion de
        arquitectura. La realidad es \rt{template\_artifacts.lifecycle\_state} con tres valores.

  \item \textbf{No hay tabla de auditoria transversal.} Solo bitacoras especificas por dominio.

  \item \textbf{\rt{docker/docs/Dockerfile} es una imagen huerfana}: no la referencia ningun
        compose ni el workflow.

  \item \textbf{\rt{docker/README.md} documenta puertos de RabbitMQ en qa y prod que los
        overlays ya no publican.}

  \item \textbf{\rt{program\_router.js} y \rt{unit\_router.js} son literalmente el mismo router
        duplicado}, y \rt{notification\_router.js} duplica dos endpoints de \rt{chat\_router.js}.

  \item \textbf{\rt{backend/public/} esta vacia} pese a montarse con \rt{express.static}.

  \item \textbf{Coexisten dos configuraciones de ESLint en el frontend}:
        \rt{eslint.config.cjs} (flat config, la efectiva con ESLint 10) y \rt{.eslintrc.js}
        (legacy, residual e inerte).

  \item \textbf{No hay \rt{.env} dentro de \rt{frontend/} ni de \rt{backend/}}: viven en
        \rt{docker/}, y el modelo de referencia es \rt{docker/.env\_model}.
\end{enumerate}

\section{Usuarios y accesos de referencia}

Los crea el \textbf{bootstrap} (\rt{/setup} $\rightarrow$ ``usar datos de ejemplo''). \textbf{No
hay ningun seed SQL alternativo.}

\begin{center}
\footnotesize
\begin{tabularx}{\textwidth}{@{}P{2.4cm}P{3.4cm}P{3.0cm}L@{}}
\toprule
\textbf{Perfil} & \textbf{Cedula} & \textbf{Contrasena} & \textbf{Notas} \\
\midrule
admin   & \rt{1234567890} & \rt{Demo1234!}   & bloqueado en el espacio de usuario \\
gestor  & \rt{0987654321} & \rt{Gestor1234!} & de momento tiene tambien rol de usuario \\
usuario & \rt{1122334455} & \rt{Demo1234!}   & \\
\bottomrule
\end{tabularx}
\end{center}

\begin{aviso}[El admin no puede probar el dossier ni las firmas]
El router bloquea el espacio de usuario para el administrador con
\rt{meta: \{ blockedForAdmin: true \}} (el guard lo redirige a \rt{/admin}): \rt{/home},
\rt{/home/documentos}, \rt{/home/firmas} y \rt{/perfil} con todas sus hijas, porque vue-router
hereda el \rt{meta}. \textbf{Para probar el dossier o las firmas hay que entrar como gestor o
como usuario.}
\end{aviso}

Puertos de \rt{dev}: proxy HTTP \textbf{8088} y HTTPS \textbf{8443} (la API bajo
\rt{/api/deasy/v1}); el puerto directo del backend es \textbf{3030}. Credenciales de SonarQube:
usuario \rt{admin}, contrasena \rt{Demo1234!Demo} --- pero \textbf{la API no acepta esa
contrasena por basic auth}: todo va por token.

% =====================================================================
\chapter{Por donde empezar a leer}
% =====================================================================

Ruta de lectura sugerida, de menor a mayor dificultad:

\begin{center}
\footnotesize
\begin{tabularx}{\textwidth}{@{}P{0.8cm}P{6.4cm}L@{}}
\toprule
\textbf{\#} & \textbf{Fichero} & \textbf{Por que} \\
\midrule
1  & \rt{CLAUDE.md} & Las reglas del proyecto, y es lo mas actualizado \\
2  & \rt{backend/index.js} & El arranque completo en unas 200 lineas \\
3  & \rt{backend/config/apiPaths.js} & El mapa de URLs de toda la API \\
4  & \rt{backend/services/documents/DocumentWorkflowResetService.js} & \textbf{El estilo objetivo}: una sola responsabilidad, se lee de una sentada \\
5  & \rt{backend/config/rbacCatalog.js} & Entiendes el modelo de permisos entero \\
6  & \rt{backend/database/postgres\_schema.sql} & Los comentarios explican el porque de cada decision \\
7  & \rt{docs/arquitecturas/modelo-emision-entregables.md} & La fuente de verdad de \rt{item\_mode} \\
8  & \rt{docs/planes/referencia/calidad-y-medicion.md} & Que esta mal, que se esta arreglando, y \textbf{que NO hay que tocar} \\
9  & \rt{frontend/src/core/router/index.js} & El mapa del frontend \\
10 & \rt{signer/app.py} & Dejalo para el final: es denso \\
\bottomrule
\end{tabularx}
\end{center}

\section{Los tres documentos maestros}

\begin{itemize}[leftmargin=1.4em]
  \item \rt{docs/planes/referencia/calidad-y-medicion.md} --- el \textbf{documento maestro} de
        deuda tecnica y complejidad: mapa de fases con su estado, linea base de SonarQube,
        ranking de ficheros y funciones, y la lista de \textbf{lo que NO hay que tocar}
        (\rt{sqlTables.js} y los falsos positivos entran ahi). \textbf{Leelo antes de proponer un
        refactor.}
  \item \rt{docs/planes/referencia/cobertura.md} --- la cobertura. Y lo primero que dice: el gate
        \textbf{no pide 80\,\% global} (eso seria trabajo de anos), pide \textbf{80\,\% de lo
        nuevo}.
  \item \rt{docs/planes/referencia/patrones-diseno.md} --- \textbf{cuando usar un patron de
        diseno y cuando no}. En este repositorio la complejidad se cura con \textbf{tablas y
        extraccion}, no con jerarquias; hay tres sitios donde un patron GoF si se gana el sueldo,
        y una lista de donde seria sobreingenieria.
\end{itemize}

\section{Las tres cosas que mas cuestan}

Si algo se te atraviesa, casi seguro es una de estas tres:

\begin{enumerate}[leftmargin=1.6em]
  \item \textbf{El motor de procesos} (serie $\rightarrow$ regla $\rightarrow$ flujo, y los tres
        \rt{item\_mode}). Capitulo 5 de este documento.
  \item \textbf{El flujo de firma completo}, desde el clic hasta el PDF firmado en MinIO, pasando
        por la cola. Capitulos 4 y 7.
  \item \textbf{El panel de administracion generico} (\rt{sqlTables.js} +
        \rt{SqlAdminService.js} + \rt{tableHooks.js}, y su contraparte
        \rt{AdminTableManager.vue}). Capitulos 4 y 6.
\end{enumerate}

\vfill
\begin{center}
\rule{0.4\textwidth}{0.4pt}\\[6pt]
\small Documento generado a partir del analisis directo del repositorio.\\
Las cifras (lineas, conteos de tests, tamanos de fichero) son fotos del momento y cambian.
\end{center}

\end{document}
